% concatenated from initial2final_arxiv1.tex
\documentclass[11pt,a4paper]{amsart}

%%%%%%%%%%%%%%%%%%%%%%%%%%%%%%%%%%%%
\usepackage{amssymb}
\usepackage[utf8]{inputenc}
\usepackage[english]{babel}
\usepackage{amsfonts,mathtools}
\usepackage{amsmath}
\usepackage{comment}
\usepackage{transparent}
\usepackage{graphicx}
\usepackage{xcolor}
\usepackage{soul}
\usepackage{tikz}
\usetikzlibrary{calc}
\setstcolor{teal}
\usepackage{todonotes}

%%%%%%%%%%%%%%%%%%%%%%%%%%% commenting
\newcommand{\prp}[1]{{\color{purple}#1}}
\newcommand{\marnote}[1]{\marginpar{\parbox[b]{0.5\marginparwidth}{\footnotesize #1}}} %margin notes


\DeclareFontFamily{U}{mathx}{\hyphenchar\font45}
\DeclareFontShape{U}{mathx}{m}{n}{
      <5> <6> <7> <8> <9> <10>
      <10.95> <12> <14.4> <17.28> <20.74> <24.88>
      mathx10
      }{}
\DeclareSymbolFont{mathx}{U}{mathx}{m}{n}
\DeclareFontSubstitution{U}{mathx}{m}{n}
\DeclareMathAccent{\widecheck}{0}{mathx}{"71}
\DeclareMathAccent{\wideparen}{0}{mathx}{"75}

\def\cs#1{\texttt{\char`\\#1}}


\usepackage[bookmarksopen,bookmarksdepth=3,colorlinks,citecolor=red,pagebackref,hypertexnames=false]{hyperref}
\usepackage[nobysame,non-sorted-cites, initials]{amsrefs}


\oddsidemargin=-.0cm
\evensidemargin=-.0cm
\textwidth=16cm
\textheight=22cm
\topmargin=0cm


\usepackage{amsthm}
\usepackage{color}
\usepackage{enumitem}
\usepackage[nameinlink]{cleveref}

\renewcommand{\eprint}[1]{\href{http://arxiv.org/abs/#1}{#1}}

%%%%%%%%%%%%%%%%%%%%%%%%%%%%%%%%%%%%
% \newtheorem{main-theorem}{Theorem}
\newtheorem{proposition}{Proposition}[section]
\newtheorem{theorem}[proposition]{Theorem}
\newtheorem{claim}[proposition]{Claim}
\newtheorem{corollary}[proposition]{Corollary}
\newtheorem{lemma}[proposition]{Lemma}
\theoremstyle{remark}
\newtheorem{remark}[proposition]{Remark}
\newtheorem{example}{Example}[section]
\theoremstyle{definition}
\newtheorem{definition}[proposition]{Definition}
\newtheorem*{acknowledgements}{Acknowledgements}

%%%%%%%%%%%%%%%%%%%%%%%%%%%%%%%%%%%%
\DeclareMathOperator{\supp}{supp}
\DeclareMathOperator{\inte}{int}
\DeclareMathOperator{\sign}{sign}
\DeclareMathOperator{\graph}{graph}

\newcommand{\R}{\mathbb{R}}
\newcommand{\C}{\mathbb{C}}
\newcommand{\Z}{\mathbb{Z}}
\newcommand{\N}{\mathbb{N}}
\newcommand{\Sph}{\mathbb{S}}
\newcommand{\dd}{\mathrm{d}}
\newcommand{\tm}{\mathrm{t}}
\newcommand{\x}{\mathrm{x}}
\newcommand{\Id}{\mathrm{Id}}
\newcommand{\Orth}{\mathrm{O}}
\newcommand{\T}{\mathsf{T}}


\newcommand{\X}{\mathrm{X}}


\def\loc{\mathop\mathrm{loc}}
\def\ve{\varepsilon}
\def\lec{\lesssim}
\def\vp{\varphi}
\def\wh{\widehat}
\def\ind{\mathbf 1}
\def\cal{\mathcal}
\def\ovl{\overline}
\def\cd{\centerdot}
\def\wt{\widetilde}
\def\d{\partial}




%%%%%%%%%%%%%%%%%%%% title %%%%%%%%%%%%%%%%%%%%
\title[Initial-to-final-state inverse problem]{The initial-to-final-state inverse problem with critically-singular potentials}


%%%%%%%%%%%%%%%%%%%% authors %%%%%%%%%%%%%%%%%%%%
\author[M. Cañizares]{Manuel Cañizares}
\address[M.\ Cañizares]{Johann Radon Institute for Computational and Applied Mathematics (RICAM)\\ Altenbergstr. 69, 4040 Linz, Austria.}
\email{\href{mailto:manuel.canizares@ricam.oeaw.ac.at}{\textrm{manuel.canizares@ricam.oeaw.ac.at}}}


\author[P. Caro]{Pedro Caro}
\address[P. Caro]{
Basque Center for Applied Mathematics and Ikerbasque (Basque Foundation for Science) Bilbao, Spain}
\email{\href{mailto:pcaro@bcamath.org}{\textrm{pcaro@bcamath.org}}}


\author[I. Parissis]{Ioannis Parissis}
\address[I.\ Parissis]{Departamento de Matem\'aticas, Universidad del Pa\'is Vasco, Aptdo. 644, 48080 Bilbao, Spain and Ikerbasque, Basque Foundation for Science, Bilbao, Spain}
\email{\href{mailto:ioannis.parissis@ehu.eus}{\textrm{ioannis.parissis@ehu.eus}}}



\author[T. Zacharopoulos]{Thanasis Zacharopoulos}
\address[T.\ Zacharopoulos]{Department of Mathematics, Aarhus University, NY Munkegade 118, 8000 Aarhus C, Denmark}
\email{\href{mailto:thanzacharop@math.au.dk}{\textrm{thanzacharop@math.au.dk}}}


%%%%%%%%%%%%%%%%%%%%%%%%%%%%%%%%%%%%
\begin{document}

%%%%%%%%%%%%%%%%%%%%%%%%%%%%%% THEOREM THEOREM THEOREM
\begin{abstract} The Schr\"odinger equation in high dimensions describes the evolution of a quantum system. Assume that we are given the evolution map sending each initial state $f\in L^2(\R^n)$ of the system to the corresponding final state at a fixed time $T$. The main question we address in this paper is whether this \emph{initial-to-final-state} map uniquely determines the Hamiltonian $-\Delta+V$ that generates the evolution. We restrict attention to \emph{time-independent} potentials $V$ and show that uniqueness holds provided $V \in L^1(\R^n)\cap L^q(\R^n)$, with $q>1$ if $n=2$ or $q\geq n/2$ if $n\geq 3$. This should be compared with the results of Caro and Ruiz, who proved that in the time-dependent case, uniqueness holds under the stronger assumption that the potential exhibits super-exponential decay at infinity, for both bounded and unbounded potentials.

This paper extends earlier work of the same authors, where uniqueness was obtained for bounded time-independent potentials with polynomial decay at infinity. Here we only require $L^1$-type decay at infinity and allow for $L^q$-type singularities. We reach this improvement by providing a refinement of the Kenig--Ruiz--Sogge resolvent estimate, which replaces the classical Agmon--H\"ormander estimates used previously. Crucially, the time-independent setting allows us to avoid the use of complex geometrical optics solutions and thereby dispense with strong decay assumptions at infinity.
\end{abstract}
%%%%%%%%%%%%%%%%%%%%%%%%%%%%%% THEOREM THEOREM THEOREM


\date{\today}

\subjclass[2020]{Primary 35R30; Secondary 35J10, 81U40.}
\keywords{inverse problems, Schr\"odinger equation, time-independent potentials,
initial-to-final-state map, uniqueness, resolvent estimates}

\maketitle

%%%%%%%%%%%%%%%%%%%%%%%%%%%%%% SECTION SECTION SECTION
\section{Introduction}
In this paper, we study an inverse problem for the Schr\"odinger equation in which the available data consist of the map that sends any initial state $f$ at time $t=0$ to the solution at a fixed final time $t=T$.

To make this precise, let us consider the initial-value problem for the Schr\"odinger equation with a time-independent potential
\begin{equation}\label{eq:Schrodinger}
\begin{cases}
i \partial_t u = -\Delta u + V u & \text{for } (t,x)\in (0, T)\times \R^n \eqqcolon \Sigma,
\vspace{.6em}
\\
u(0,x) = f(x) & \text{for } x\in \R^n.
\end{cases}
\end{equation}
We assume throughout the paper that $V=V(x)\in L^q(\R^n)$, where $q\ge n/2$ if $n\ge 3$ or $q>1$ if $n=2$. Then, by \cites{zbMATH02204588,zbMATH00179225}, this direct problem is well posed, and for every $f\in L^2(\R^n)$ there exists a unique solution
\[
u\in C\left([0,T];L^2(\R^n)\right).
\]
The evolution associated with \eqref{eq:Schrodinger} therefore defines a bounded linear operator
\[
\mathcal U : f\in L^2(\R^n)\mapsto u\in C\left([0,T];L^2(\R^n)\right).
\]
Consequently, for any fixed time $t\in[0,T]$, the operator
\[
\mathcal U_t : f \in L^2(\R^n) \mapsto u(t,\cd)\in L^2(\R^n)
\]
is also bounded, uniformly in $t$. Solutions of the form $u=\mathcal U f$, with $f\in L^2(\R^n)$, will be referred to as \emph{physical solutions}, while we call $\mathcal U_T$ the \emph{initial-to-final-state map}.

The main question addressed in this paper is whether the initial-to-final-state map $\mathcal U_T$ uniquely determines the Hamiltonian $-\Delta+V$. This inverse problem was first studied for \emph{time-dependent} potentials in \cite{zbMATH07801151}. There, the authors show that if the potentials $V_1,V_2 \in L^1((0,T);L^\infty(\R^n))$ satisfy a \emph{super-exponential decay} condition at infinity, and if $\mathcal U_T^j$ denotes the initial-to-final-state map associated with $-\Delta+V_j$, then there holds:
\[
\mathcal U_T^1=\mathcal U_T^2 \quad\Longrightarrow\quad V_1=V_2.
\]
More recently, this uniqueness result was extended in \cite{caro2025initialtofinalstateinverseproblemunbounded} to unbounded time-dependent potentials that are allowed to exhibit local $L^q$-type singularities, but still requiring the super-exponential decay assumption at infinity.

The case of time-independent potentials was previously considered in \cite{zbMATH08122191}, where uniqueness was established under comparatively weaker decay assumptions than in the time-dependent setting, namely assuming only super-linear decay at infinity. The purpose of the present paper is to relax these assumptions even further. Specifically, we prove uniqueness for time-independent potentials that may exhibit singularities of $L^q$-type in sets of finite measure and satisfy only $L^1$-integrability in sets of infinite measure; in the time-independent setting, this represents a substantial improvement over the decay and integrability assumptions in  \cite{zbMATH08122191}, \cite{zbMATH07801151}, and \cite{caro2025initialtofinalstateinverseproblemunbounded}.

%%%%%%%%%%%%%%%%%%%%%%%%%%%%%% THEOREM THEOREM THEOREM
\begin{theorem}\sl \label{th:Lq} Let $V_1,V_2\in L^1(\R^n)\cap L^q(\R^n)$ be time-independent potentials, where $q>1$ if $n=2$ and $q\ge n/2$ if $n\ge 3$. Let $\mathcal U_T^1$ and $\mathcal U_T^2$ denote the corresponding initial-to-final-state maps. Then
\[
\mathcal U_T^1=\mathcal U_T^2 \quad\Longrightarrow\quad V_1=V_2.
\]
\end{theorem}
%%%%%%%%%%%%%%%%%%%%%%%%%%%%%% THEOREM THEOREM THEOREM

\medskip
\noindent\textbf{Outline of the proof.}
The proof of Theorem~\ref{th:Lq} proceeds by extracting information on the difference $V_1-V_2$ from $\mathcal U_T^1 = \mathcal U_T^2$, by testing it against suitable families of solutions.

The first step is to show that the equality $\mathcal U_T^1=\mathcal U_T^2$ yields an \emph{Alessandrini-type orthogonality relation}, in the sense of \cite{Aless}, of the form
\begin{equation}\label{eq:orthogonality_intro}
\int_{\Sigma} (V_1-V_2)\,u_1\,\overline{v_2}=0,
\end{equation}
valid for pairs of solutions associated with the potentials $V_1$ and $\ovl{V_2}$. This identity is initially available only for \emph{physical solutions}, namely solutions belonging to $C([0,T];L^2(\R^n))$. 

At this stage, our first obstruction is that in order to recover pointwise information on $V_1-V_2$, we need to test \eqref{eq:orthogonality_intro} against special \emph{time-harmonic} solutions. Such solutions can be constructed as perturbations of time-harmonic solutions of the  free Schr\"odinger equation; we refer to such solutions as \emph{stationary states} associated with the given potential. However, these stationary states are not in $C([0, T]; L^2(\R^n))$, and using them as test functions requires a substantial extension of the Alessandrini-type orthogonality relation beyond physical solutions.

More precisely, stationary states are obtained by inverting a resolvent-type operator and constructing correction terms via Neumann series. Concretely, they are of the form
\[
u(t,x)=e^{-i|\kappa|^2 t}\left(e^{-i\kappa\cdot x}+w^{\mathrm{cor}}(x)\right),
\]
where the leading term \(e^{-i|\kappa|^2 t}e^{-i\kappa\cdot x}\) is a time-harmonic solution of the free Schr\"odinger equation, and the correction \(w^{\mathrm{cor}}\) accounts for the presence of the potential \(V\).

The construction of \(w^{\mathrm{cor}}\) reduces to inverting an operator of the form \(\mathrm{Id}-P_\lambda\circ V\), where \(P_\lambda\) denotes a solution operator for the Helmholtz equation \((\Delta+\lambda^2)u=f\). This inversion is carried out on suitable function spaces, chosen so that the operator \(P_\lambda\circ V\) is small, in a suitable sense.

In following the proof strategy outlined above, we encounter a second important obstruction. In the non-endpoint regime $V\in L^q$ with $q>n/2$, decay in the energy parameter $\lambda$ follows from the classical Kenig--Ruiz--Sogge resolvent estimate. The latter ensures that $P_\lambda\circ V$ is small on the relevant function spaces for large $\lambda$. At the critical endpoint $q=n/2$, however, this estimate no longer yields decay in $\lambda$, and the smallness of $P_\lambda\circ V$ cannot be obtained from the standard resolvent bound alone.

To address this issue, we introduce a new scale of Banach spaces that allows us to recover a decaying factor in the energy parameter, also at the endpoint. Within this framework, we establish an improved form of the Kenig--Ruiz--Sogge resolvent estimate, sufficient to construct stationary states for critical potentials despite the absence of quantitative decay at the endpoint. A similar idea has appeared in \cite{CaroGarcia} for a related scattering problems with critically singular potentials; see also \cite{zbMATH07867333}. When inserted into the extended Alessandrini-type orthogonality relation, these stationary states allow us to isolate the Fourier phase and hence recover the Fourier transform of $V_1-V_2$, proving uniqueness.

Throughout the argument, the assumption of a stationary potential plays a crucial role. Indeed, in the time-dependent setting, the uniqueness results in \cites{zbMATH07801151,caro2025initialtofinalstateinverseproblemunbounded} rely on constructing complex geometrical optics solutions whose leading terms are of the form
\[
(t,x) \longmapsto e^{it|\kappa|^2} e^{\kappa\cdot x}, \qquad \kappa \in \R^n,
\]
supplemented by correction terms depending on the potential. These complex exponentials grow in certain directions.  Using such solutions in an orthogonality relation of the form \eqref{eq:orthogonality_intro} and controlling the corresponding correction terms naturally leads to assuming super-exponential decay of the potential at infinity.

However, in the time-independent case considered here, we work with time-harmonic solutions whose main term is given in the form
\[
(t,x) \longmapsto e^{-i|\kappa|^2 t} e^{-i\kappa\cdot x}, \qquad \kappa \in \R^n,
\]
ignoring again the perturbative terms correcting for the potential. Now the exponents are imaginary, so the leading term is purely oscillatory and unimodular. This allows us to restrict attention to a more regular class of solutions, which however is still sufficient to establish uniqueness in the case of time-independent potentials.

Inverse problems associated with the dynamical Schrödinger equation have been extensively studied; see, for example, \cites{zbMATH01886353,zbMATH06733553,zbMATH05549395,zbMATH05655673,zbMATH05839237,zbMATH06864429}. Many of these works have also considered time-dependent Hamiltonians \cites{zbMATH06769718,zbMATH06516179,zbMATH05379127,zbMATH07033617,zbMATH07242805}. A common feature of these investigations is the use of a dynamical Dirichlet-to-Neumann map, recorded on the boundary of a domain that contains the non-constant portions of the Hamiltonian. 

This formulation stands in contrast to the one presented in our work. Here, the variable part of the Hamiltonian is not localized to a bounded region but is possibly present throughout the whole space. Moreover, our inverse problem is formulated with a distinct data requirement: knowledge only of the system's initial state and its corresponding state at a final time.

The rest of the paper is organized as follows. Section~\ref{sec:aux} collects background material and auxiliary analytic tools, including well–posedness results and resolvent estimates. In Section~\ref{sec:timeindependent_solutions} we construct stationary states for the Schr\"odinger equation with time-independent potentials, treating both the non-endpoint case $q>n/2$ and the endpoint case $q=n/2$ in dimensions $n\ge3$. Section~\ref{sec:orthorelation} is devoted to extending the Alessandrini-type orthogonality relation beyond physical solutions, and in particular to stationary-state solutions. Finally, in Section~\ref{sec:unbounded_endpoint} we prove Theorem~\ref{th:Lq} by applying the extended orthogonality relation to the stationary states constructed earlier.


%%%%%%%%%%%%%%%%%%%%%%%%%%%%%% SECTION SECTION SECTION
\section{Preliminaries: Function spaces, Resolvent and Strichartz  estimates} \label{sec:aux} In this section we collect the basic functional-analytic tools and estimates used throughout the paper.

Throughout the paper, we use the Fourier transform on $\R^n$ defined for $f\in L^1(\R^n)$ by
\[
\widehat f(\xi) \coloneqq \frac{1}{(2\pi)^{n/2}} \int_{\R^n} f(x) e^{-i x\cdot \xi}\, \dd x,
\qquad \xi\in\R^n,
\]
and extended in the usual way to tempered distributions, in particular to $L^p(\R^n)$ for $1<p\le2$.

For $p,r\in[1,\infty]$ and an open time interval $I\subset\R$, we use the Banach spaces $L^r\left (I;L^p(\R^n)\right )$ and $C\left (\overline{I};L^p(\R^n)\right )$, equipped with the norms
\begin{equation}\label{eq:Bvalued_norms}
\|u\|_{L^r\left(I;L^p(\R^n)\right)} \coloneqq
\left( \int_I \|u(t,\cd)\|_{L^p(\R^n)}^r\,\dd t \right)^{1/r}, \qquad \|u\|_{C\left(\overline{I};L^p(\R^n)\right)} \coloneqq \sup_{t\in I}\|u(t,\cd)\|_{L^p(\R^n)}.
\end{equation}


We will also use intersections of mixed-norm spaces of the form
\[
L^{r_1}\left (I;L^{p_1}(\R^n)\right )\cap L^{r_2}\left(I;L^{p_2}(\R^n)\right), \textnormal{ and } L^{p_1}(\R^n)\cap L^{p_2}(\R^n)
\]
endowed with the norm
\[
\max\left( \|u\|_{L^{r_1}\left(I;L^{p_1}(\R^n)\right)}, \|u\|_{L^{r_2}\left(I;L^{p_2}(\R^n)\right)}\right) , \textnormal{ and } \max\left( \|u\|_{L^{p_1}(\R^n)}, \|u\|_{L^{p_2}(\R^n)}\right)
\]
respectively.


%%%%%%%%%%%%%%%%%%%%%%%%%%%%%% SECTION SECTION SECTION
\subsection{The Kenig--Ruiz--Sogge resolvent estimate}\label{sec:KRS} For $\lambda>0$, we consider the Helmholtz equation
\begin{equation}\label{eq:Helm}
(\Delta+\lambda^2)u=f \qquad \text{in } \R^n.
\end{equation}
For $f\in\mathcal S(\R^n)$, we choose the associated solution operator defined by
\begin{equation}\label{eq:soloper}
P_\lambda f(x) \coloneqq \frac{1}{(2\pi)^{n/2}} \,\mathrm{p.v.}\!\int_{\R^n} \frac{e^{ix\cdot\xi}}{\lambda^2-|\xi|^2} \widehat f(\xi)\,\dd\xi, \qquad x\in\R^n.
\end{equation}
Then $u\coloneqq P_\lambda f$ is a distributional solution of \eqref{eq:Helm}.

We denote by $q_n$ the Lebesgue exponent arising from the Tomas--Stein extension theorem, \cites{stein,Tomas}, defined by
\begin{equation}\label{eq:extLq}
\frac{1}{q_n} \coloneqq \frac{1}{2}-\frac{1}{n+1}=\frac{n-1}{2(n+1)}, \qquad n\geq 2.
\end{equation}
This exponent appears naturally in the analysis of the Helmholtz equation \eqref{eq:Helm} below, and in particular in the $L^p$-bounds for its solution operator.

Let $p_n$ denote the Sobolev exponent for the embedding $\dot H^1(\R^n)\hookrightarrow L^{p_n}(\R^n)$, namely
\[
\frac{1}{p_n}=\frac12-\frac1n=\frac{n-2}{2n}, \qquad n\geq 3.
\]
The following quantitative estimate for $P_\lambda$ is obtained from the work of Kenig, Ruiz, and Sogge \cite{zbMATH04050093} by rescaling:
\begin{equation}\label{eq:KRS}
\|P_\lambda f\|_{L^p(\R^n)} \lesssim \frac{1}{\lambda^{2n(\frac1p-\frac1{p_n})}} \|f\|_{L^{p'}(\R^n)} \quad\text{for}\quad
\begin{cases} 
q_n\leq p\le p_n, & n\geq3, 
\vspace{.6em}
\\ 
q_n\leq p<\infty, & n=2,
\end{cases}
\end{equation}
where the implicit constant depends only on $p$ and $n$. From here on, $p'$ denotes the H\"older conjugate exponent of $p$.



%%%%%%%%%%%%%%%%%%%%%%%%%%%%%% SECTION SECTION SECTION
\subsection{Strichartz pairs and well-posedness}\label{subsec: Stri_WP} We assume $n\geq 2$ throughout the paper. We will say that the pair $(r, p)$ is a \emph{Strichartz pair} if 
\[
\frac{2}{r} + \frac{n}{p} = \frac{n}{2}, \quad \textup{with}\quad r, p \in [2, \infty] \quad  \mathrm{and} \quad (n, r, p) \neq (2, 2, \infty).
\]
Note that, in dimension $n=2$, the endpoint pair $(2,\infty)$ is not admissible.

We consider the Schr\"odinger equation in $\Sigma\coloneqq [0,T]\times \R^n$ with a time-independent potential $V=V(x)$ and a force term $F=F(t,x)$ as follows
\begin{equation}\label{eq:schfull}
\begin{cases}
(i\partial_t +\Delta)u-Vu=F,\quad &(t,x)\in \Sigma,
\vspace{.6em}
\\
u(0,\cd)=f, \quad &\textrm{in}\quad \R^n.
\end{cases}
\end{equation}
Assume that the potential $V$ is time-independent and satisfies $V\in L^q(\R^n)$, where $q\geq n/2$ if $n\ge3$ and $q>1$ if $n=2$. Let $(r,p)$ and $(\tilde r,\tilde p)$ be Strichartz pairs satisfying $\frac{1}{q}+\frac{2}{p}=1$. Then, for every $f\in L^2(\R^n)$ and $F\in L^{\tilde r'}((0,T);L^{\tilde p'}(\R^n))$, the initial value problem \eqref{eq:schfull} admits a unique solution
\[
u\in C\left([0,T];L^2(\R^n)\right)\cap L^r\left((0,T);L^p(\R^n)\right);
\]
see Figure~\ref{fig:str}. This follows from the existence and uniqueness results for critical potentials proved by Ruiz--Vega \cite{zbMATH00179225} and Ionescu--Kenig \cite{zbMATH02204588}.



Throughout the paper, the integrability exponents for the potential and the solutions are linked by the H\"older relation
\[
\frac{1}{q}+\frac{1}{p}+\frac{1}{p}=1.
\]
Under this assumption, the following bounds hold:
\begin{equation}\label{eq:holder}
\left|\int_{\R^n} V u v\right| \leq \|V\|_{L^q(\R^n)} \|u\|_{L^p(\R^n)}\|v\|_{L^p(\R^n)},\qquad \|Vu\|_{L^{p'}(\R^n)}\leq \|V\|_{L^q(\R^n)}\|u\|_{L^p(\R^n)}.
\end{equation}
If, in addition, $r\geq 2$, then for any time interval $(0,T)$ one also has
\begin{equation}\label{eq:holder1}
\|Vu\|_{L^{r'} \left((0,T);L^{p'}(\R^n)\right)} \leq T^{\frac{1}{r'}-\frac{1}{r}} \|V\|_{L^q(\R^n)} \|u\|_{L^r \left((0,T);L^p(\R^n)\right)}.
\end{equation}

%%%%%%%%%%%%%%%%%%%%%%%%%%%%%% REMARK REMARK REMARK
\begin{remark}\label{rmrk:VuinLp'} Let $V\in L^q(\R^n)$ be time independent, with $q>1$ if $n=2$ or $q\geq n/2$ if $n\geq3$, and let $(r,p)$ be a Strichartz pair satisfying $\frac1q+\frac2p=1$. If $u\in C([0,T];L^2(\R^n))\cap L^r((0,T);L^p(\R^n))$ solves \eqref{eq:schfull} with forcing term $F\in L^{r'}((0,T);L^{p'}(\R^n))$, then \eqref{eq:holder1} implies
\[
(i\partial_t+\Delta)u = F+Vu \in L^{r'}\big((0,T);L^{p'}(\R^n)\big).
\]
In particular, this holds for the endpoint Strichartz pair $(r_n,p_n)=(2,\frac{2n}{n-2})$ when $n\ge3$. This property will be used repeatedly in Section~\ref{sec:orthorelation} and uses the fact that $V$, or the inverse problem itself, is only considered throughout a time interval of finite length, $T<\infty$.
\end{remark}
%%%%%%%%%%%%%%%%%%%%%%%%%%%%%% REMARK REMARK REMARK


%%%%%%%%%%%%%%%%%%%%%%%%%%%%%% FIGURE FIGURE FIGURE
\begin{figure}[htb]
\centering
\def\svgwidth{300pt}
\input{drawing.pdf_tex}
\vspace{-5em}
\caption[Strichartz line]{\small The exponents for the space of potentials $L^q(\R^n)$, the Strichartz pairs $(r,p)$, and the Stein-Tomas extension estimate $L^2(\mathbb S^{n-1})\to L^{p}(\R^n)$ for $p\geq q_n$. The Kenig--Ruiz--Sogge estimate is valid in the range $p\in [q_2,\infty)$ when $n=2$ and $p\in[q_n,p_n]$ when $n\geq 3$.} \label{fig:str}
\end{figure}
%%%%%%%%%%%%%%%%%%%%%%%%%%%%%% FIGURE FIGURE FIGURE

%%%%%%%%%%%%%%%%%%%%%%%%%%%%%% REMARK REMARK REMARK
\begin{remark}\label{rmrk:numer} A brief comment on the role of the exponents is in order. The Kenig--Ruiz--Sogge estimate \eqref{eq:KRS} holds in the range
\[
q_n \le p \leq p_n \quad \text{if } n\geq3, \qquad q_n \le p < \infty \quad \text{if } n=2.
\]
Imposing the H\"older relation $\frac{2}{p}+\frac{1}{q}=1$, this corresponds to
\[
\frac{n}{2} \leq q \leq \frac{n+1}{2} \quad \text{if } n\geq3, \qquad 1< q \leq \frac{3}{2} \quad \text{if } n=2.
\]
In dimensions $n\ge3$, the endpoint $q=n/2$ corresponds to $p=p_n$, and in this case \eqref{eq:KRS} no longer yields decay in the parameter $\lambda$, which necessitates a separate treatment based on a refinement of \eqref{eq:KRS}, developed in Section~\ref{sec:timeindependent_solutions} below. By contrast, the upper endpoint $q=(n+1)/2$, corresponding to $p=q_n$, does not pose any difficulty in our setting. Indeed, since throughout the paper we assume $V\in L^1(\R^n)\cap L^q(\R^n)$, any potential with $q>(n+1)/2$ automatically belongs to $L^1(\R^n)\cap L^{(n+1)/2}(\R^n)$, with the norm bound
\[
\|V\|_{L^1(\R^n)\cap L^{\frac{n+1}{2}}(\R^n)} \leq \|V\|_{L^1(\R^n)\cap L^q(\R^n)}.
\]
Thus, without loss of generality, we may always restrict attention to the range $q\leq (n+1)/2$.
\end{remark}
%%%%%%%%%%%%%%%%%%%%%%%%%%%%%% REMARK REMARK REMARK

%%%%%%%%%%%%%%%%%%%%%%%%%%%%%% SECTION SECTION SECTION 
\section{Solutions of the time-independent Schr\"odinger equation}\label{sec:timeindependent_solutions}
In this section we construct the special solutions, called stationary states, that will be used in the orthogonality relation~\eqref{eq:orthogonality_intro} to prove Theorem~\ref{th:Lq} in Section~\ref{sec:unbounded_endpoint}. Specifically, we consider functions $\psi$ solving the Schr\"odinger equation
\begin{equation}\label{sch_timeind}
(i\partial_t +\Delta-V)\psi =0 \qquad\mathrm{in}\quad \R\times \R^n,\quad n\geq 2,
\end{equation}
of the form
\begin{equation}\label{eq:steadysol}
\psi_V ^{\lambda,\omega}(t,x)\coloneqq e^{-i\lambda^2 t}(e^{-i\lambda \omega\cdot x}+w^{\mathrm{cor}} (x))\eqqcolon e^{-i\lambda^2 t}\big(w^{(0)}(x)+w^{\mathrm{cor}}(x)\big),\qquad  (t,x)\in\R\times \R^n,
\end{equation}
where $\omega\in\mathbb S^{n-1}$ and  $\lambda>0$. Of course, the functions $w^{(0)},w^{\mathrm{cor}}$ also depend on $\lambda,\omega$ and $V$ but we suppress this dependence in order to simplify the notation. We will also use the notation
\[
\psi_V ^{\lambda,\omega}(t,x)=e^{-i\lambda^2 t}w(x),\qquad w(x)\coloneqq w^{(0)}(x)+w^{\mathrm{cor}}(x),\qquad (t,x)\in\R\times \R^n.
\]

In order for $\psi$ to solve \eqref{sch_timeind} we check that $w^{\mathrm{cor}}$ has to satisfy
\[
\left(\Delta+\lambda^2-V\right)(w^{(0)}+w^{\mathrm{cor}})=0\qquad\mathrm{in}\quad\R^n,
\]
or equivalently
\[
\left(\Delta+\lambda^2\right)w^{\mathrm{cor}}=Vw^{(0)}+Vw^{\mathrm{cor}}\qquad\mathrm{in}\quad\R^n.
\]
By the discussion in \S\ref{sec:KRS} and with the notation therein we note that the above partial differential equation will be satisfied if
\[
w^{\mathrm{cor}}=P_\lambda(Vw^{(0)}+Vw^{\mathrm{cor}}),
\]
where $P_\lambda$ is the solution operator defined in \eqref{eq:soloper}. Such $w^{\mathrm{cor}}$ would thus satisfy
\[
\left(\mathrm{Id}-P_\lambda \circ V\right)w^{\mathrm{cor}}
= P_\lambda(Vw^{(0)}).
\]
We will therefore be able to construct $w^{\mathrm{cor}}$ once we manage to invert the operator $\mathrm{Id}-P_\lambda\circ V$ in $\mathcal L(X^*)$ for a suitable Banach space $X^*$.

Throughout this section we assume that
\[
V\in L^1(\R^n)\cap L^q(\R^n),\qquad  \frac n2 \leq q\leq \frac {n+1}{2};
\]
see Remark~\ref{rmrk:numer} concerning the range of $q$. Since $w^{(0)}$ is unimodular, we have $Vw^{(0)}\in L^1(\R^n)\cap L^q(\R^n)$. Consequently, in order to construct $w^{\mathrm{cor}}$ we will choose a Banach space $X^*$ such that
\[
P_\lambda\left(L^1(\R^n)\cap L^q(\R^n)\right)\subseteq X^*,
\]
and for which $\mathrm{Id}-P_\lambda\circ V$ can be inverted.

We henceforth divide the discussion and the corresponding choice of $X^*$ according to whether $q>n/2$ for $n\geq 2$, corresponding to the non-endpoint case, or $q=n/2$ for $n\geq 3$, corresponding to the critical case.


%%%%%%%%%%%%%%%%%%%%%%%%%%%%%% SECTION SECTION SECTION
\subsection{Non-endpoint solutions for \texorpdfstring{$q>n/2$}{q>n/2}} Consider a \emph{time-independent} potential $V\in L^1(\R^n)\cap L^q(\R^n)$ for some $n/2<q\leq (n+1)/2$, and define  $p$ by the H\"older relation
\[
\frac{1}{q}+\frac{2}{p}=1.
\]
We begin by noting that $1<p'<q$ and thus $L^1(\R^n)\cap L^q(\R^n)\subset L^{p'}(\R^n)$. Remembering the mapping properties of $P_\lambda$ from \eqref{eq:KRS} we have that 
\[
P_\lambda\left(L^1(\R^n)\cap L^q(\R^n)\right)\subset P_\lambda\big (L^{p'}(\R^n)\big)\subset L^p(\R^n).
\]
Thus we will be able to construct the correction term $w^{\mathrm{cor}}$ in $L^p(\R^n)\eqqcolon X^*$ once we manage to show that the operator $\mathrm{Id}-P_\lambda\circ V$ is invertible in $\mathcal L(X^*)=\mathcal L(L^p(\R^n))$. This is the content of the following lemma.

%%%%%%%%%%%%%%%%%%%%%%%%%%%%%% LEMMA LEMMA LEMMA
\begin{lemma}\label{lem:I-P inverse in Lp}\sl Let $n\geq 2$ and $V \in L^q(\R^n) $ for some $n/2< q \leq (n+1)/2$. Define $p$ via $\frac{1}{q} = 1 - \frac{2}{p} $. There exists a constant $\lambda_V$ depending on $n,q$ and $\|V\|_{L^q(\R^n)}$ such that for every $ \lambda > \lambda_V $, the operator $ \Id - P_\lambda \circ V $ has a bounded inverse in $\mathcal{L}(L^p(\R^n))$. Moreover, if this inverse is denoted by $ ( \Id - P_\lambda \circ V )^{ -1 }$ we have that
\[
 \left\| ( \Id - P_\lambda \circ V )^{-1} \right\|_{\mathcal{L}(L^p(\R^n))} \leq  2 
\]
for all $\lambda \geq 2\lambda_V $.
\end{lemma}
%%%%%%%%%%%%%%%%%%%%%%%%%%%%%% LEMMA LEMMA LEMMA

%%%%%%%%%%%%%%%%%%%%%%%%%%%%%% PROOF PROOF PROOF
\begin{proof} We just note that under the assumptions on $q$ we have that $p=2q' \in[q_n,p_n)$;
hence the Kenig--Ruiz--Sogge estimate \eqref{eq:KRS} applies. Thus
\[
\left \|(P_\lambda \circ V )F\right\|_{L^p(\R^n)}\lesssim \frac{1}{\lambda^{2n\left(\frac1p-\frac{1}{p_n}\right)}} \|VF\|_{L^{p'}(\R^n)}\leq \frac{1}{\lambda^{n\left(\frac2n-\frac{1}{q}\right)}}\|V\|_{L^q(\R^n)}\|F\|_{L^p(\R^n)}
\]
by \eqref{eq:holder}, the definition of $p_n$ and the relation between $p,q$. Since $V\in  L^q(\R^n)$ it follows that $\|P_\lambda \circ V\|_{\mathcal L(L^p(\R^n))}<1$ if $\lambda>\lambda_V$ for sufficiently large $\lambda_V$ as in the statement of the lemma. The proof is now completed by a standard argument involving Neumann series.
\end{proof}
%%%%%%%%%%%%%%%%%%%%%%%%%%%%%% PROOF PROOF PROOF

Returning to the solutions in \eqref{eq:steadysol}, we can now use the invertibility of $\mathrm{Id}-P_\lambda\circ V$ to complete the construction of the correction term $w^{\mathrm{cor}}$. In the statement of the corollary below, the additional assumption $V\in L^1(\R^n)$ is used only to ensure that $Vw^{(0)}\in L^{p'}(\R^n)$, so that $P_\lambda(Vw^{(0)})$ is well defined and belongs to $L^p(\R^n)$; the invertibility of $\mathrm{Id}-P_\lambda\circ V$ on $L^p(\R^n)$ relies solely on the $L^q$ assumption.
%%%%%%%%%%%%%%%%%%%%%%%%%%%%%% COROLLARY COROLLARY COROLLARY
\begin{corollary}\label{cor:solutions}\sl Let $n\geq 2$ and consider a time-independent potential $V\in L^1(\R^n)\cap L^q(\R^n)$ for some $n/2<q\leq (n+1)/2$, and let $p$ be defined via $\frac{1}{q} = 1 - \frac{2}{p} $. For $\lambda > 2\lambda_V$ with $\lambda_V$ as in the statement of Lemma~\ref{lem:I-P inverse in Lp}, the functions
\[
w(x)\coloneqq e^{-i\lambda\omega\cdot x}+w^{\mathrm{cor}}(x)= w^{(0)}(x)+w^{\mathrm{cor}}(x),\qquad w^{\mathrm{cor}}(x)\coloneqq (\mathrm{Id}-P_\lambda\circ V)^{-1}[P_\lambda(Vw^{(0)})],
\]
with $ \omega\in\mathbb S^{n-1}$ and  $x\in\R^n$,
are weak solutions of the time-independent Schr\"odinger equation 
\[
(\Delta+\lambda^2-V)w=0\quad\mathrm{in}\quad\R^n.
\]
Additionally, for $\lambda>2\lambda_V$ we have the estimate
\[
\left \|w^{\mathrm{cor}}\right\|_{L^p(\R^n)}\lesssim \frac{1}{\lambda^{2n\left(\frac1p-\frac{1}{p_n}\right)}} \|V\|_{L^{p'}(\R^n)}\leq \frac{1}{\lambda^{2n\left(\frac1p-\frac{1}{p_n}\right)}} \|V\|_{L^1(\R^n)\cap L^q(\R^n)}.
\]
In particular, the function 
\[
\psi_V ^{\lambda,\omega}(t,x)\coloneqq e^{-i\lambda^2 t}w(x)= e^{-i\lambda^2 t}(w^{(0)}(x)+w^{\mathrm{cor}}(x)),\qquad (t,x)\in\R\times \R^n,
\]
is a weak solution of \eqref{sch_timeind}.
\end{corollary}
%%%%%%%%%%%%%%%%%%%%%%%%%%%%%% COROLLARY COROLLARY COROLLARY


%%%%%%%%%%%%%%%%%%%%%%%%%%%%%% SECTION SECTION SECTION
\subsection{Endpoint solutions in the critical case \texorpdfstring{$q=n/2,\, n\geq 3$}{q=n/2, n>=3}}\label{sec:stend} Here we assume that the potential $V\in L^1(\R^n)\cap L^{n/2}(\R^n)$. We point out that the decay in $\lambda$ in the $L^{p'}(\R^n)\to L^p(\R^n)$ estimate for $P_\lambda$ from \eqref{eq:soloper} was critical in the proof of Lemma~\ref{lem:I-P inverse in Lp} above, and thus also the subsequent Corollary~\ref{cor:solutions}. A second element used in the proof of Lemma~\ref{lem:I-P inverse in Lp} is that multiplication by $V$ is bounded from $L^p(\R^n)$ to $L^{p'}(\R^n)$, as can be read from \eqref{eq:holder} and used in the Neumann series argument.

Isolating the ingredients above, it will suffice to construct a space $X^*=X^* _\lambda \supseteq P_\lambda(L^1(\R^n)\cap L^{n/2}(\R^n))$ depending on $\lambda$, such that the following hold
\begin{align}\label{eq:modRKS}
\|P_\lambda f\|_{X^*}&\leq c(\lambda) \|f\|_{X},      
\\ \label{eq:multV}
\left|\int_{\R^n} V f g\right|&\leq \tilde c(\lambda) \|f\|_{X^*} \|g\|_{X^*},
\end{align}
with $\lim_{\lambda\to +\infty}c(\lambda)\tilde c(\lambda)=0$.

The first estimate above is a modified Kenig--Ruiz--Sogge estimate with some constant $c(\lambda)$. The second estimate  tells us that the operator of multiplication by $V$ is bounded from $X^*$ to $X$ with constant $\tilde c(\lambda)$.  Assuming these two estimates for a moment we have for any $F\in X^*$ that
\[
\left\|(P_\lambda \circ V)F\right\|_{X^*} \leq c(\lambda) \|VF\|_{X}\lesssim c(\lambda)\tilde c(\lambda)\|F\|_{X^*}
\]
and the operator norm of $P_\lambda\circ V$ becomes smaller than $1$ as $\lambda\to \infty$, allowing for the Neumann series argument.

In the remainder of this subsection we construct the space $X^*$ and prove the estimates \eqref{eq:modRKS} and \eqref{eq:multV}. To that end let us revisit the estimate of Kenig, Ruiz and Sogge, \eqref{eq:KRS}. As discussed in Remark~\ref{rmrk:numer}, the case of critical potential corresponds to $n\geq 3$ and  $q=n/2$ which happens exactly when $p=p_n$, since $\frac{2}{p}+\frac{1}{q}=1$. In this case, the estimate for the solution operator $P_\lambda$ in \eqref{eq:soloper} has no decay in $\lambda$. Since $2n(\frac{1}{q_n}-\frac{1}{p_n})=\frac{2}{n+1}$, the endpoint cases of \eqref{eq:KRS} can be written in the form
\[
      \|P_\lambda f\|_{L^{p_n}(\R^n)}\lesssim \|f\|_{L^{p_n '}(\R^n)},\qquad
     \lambda^{\frac{1}{n+1}} \|P_\lambda f\|_{L^{q_n}(\R^n)}\lesssim \frac{1}{\lambda^\frac{1}{n+1}}\|f\|_{L^{q_n '}(\R^n)}.
\]
We will construct the space $X^*$ so that it suitably combines the two estimates above. More precisely define the Banach space
\[
\X_\lambda \coloneqq L^{q_n'}(\R^n)+L^{p_n'}(\R^n) ,\qquad  \| f \|_{\X_\lambda} \coloneqq \inf \left\{ \lambda^{- \frac{1}{n+1}} \| g\|_{L^{q_n'}(\R^n)} + \| h\|_{L^{p_n'}(\R^n)} : \, f = g+h \right\}.
\]
We recall that the dual of $X_\lambda$ is isometrically isomorphic to the space
\[
X_{\lambda} ^*=  L^{q_n}(\R^n) \cap L^{p_n}(\R^n),\qquad \|g\|_{X_\lambda ^*} =\max\left\{ \lambda^{\frac{1}{n+1}} \| g \|_{L^{q_n}(\R^n)}, \| g \|_{L^{p_n}(\R^n)} \right\} .      
\]
 With $X_\lambda$ as our best candidate for the space $X$ let us first verify some easy facts. Since $1<p_n^\prime <q_n^\prime<n/2$ we have that $ L^1(\R^n)\cap L^{n/2}(\R^n) \subset L^{q_n'}(\R^n)$ and $ L^1(\R^n)\cap L^{n/2}(\R^n) \subset L^{p_n'}(\R^n)$, hence $L^1(\R^n)\cap L^{n/2}(\R^n)\subset X_\lambda$. Thus, if we manage to establish~\eqref{eq:modRKS}, then we obtain that
\[
P_\lambda \big(L^1(\R^n)\cap L^{n/2}(\R^n)\big)\subset P_\lambda(X_\lambda ) \subset X_\lambda ^*
\]
as desired. 

%%%%%%%%%%%%%%%%%%%%%%%%%%%%%% REMARK REMARK REMARK
\begin{remark}\label{rem:endpoint_norm_for_large_V} Let $\lambda>0$ and $V\in L^{n/2}(\R^n)$. Define 
\[
E_\lambda \coloneqq \left\{ x\in\R^n:\, |V(x)| > \lambda \|V\|_{L^{n/2}(\R^n)} \right\}.
\] 
Note that $V \ind_{\R^n \setminus E_\lambda} \in L^{\frac{n+1}{2}}(\R^n)$ with the estimate
\[
\left \|V \ind_{\R^n \setminus E_\lambda} \right\|_{L^{\frac{n+1}{2}}(\R^n)} \leq \lambda^{\frac{1}{n+1}} \| V \| _{L^{\frac{n}{2}}(\R^n)} .
\]
Defining
\[
\|V\|_{\lambda}\coloneqq \left \|V\ind_{E_\lambda}\right \|_{L^{n/2}(\R^n)}+\lambda^{-\frac{2}{n+1}}\left\|V \ind_{\R^n \setminus E_\lambda} \right\|_{L^{\frac{n+1}{2}}(\R^n)}
\]
and using the previous estimate, the absolute continuity of the integral of $|V|^{n/2}$ together with the fact the $|E_\lambda|\to 0$ as $\lambda\to \infty$, we conclude that $\|V\|_\lambda \to 0$ as $\lambda\to \infty$.
\end{remark} 
%%%%%%%%%%%%%%%%%%%%%%%%%%%%%% REMARK REMARK REMARK

The following lemma is the analogue of \eqref{eq:holder}, namely a H\"older-type inequality in the spirit of the desired \eqref{eq:multV}.

%%%%%%%%%%%%%%%%%%%%%%%%%%%%%% LEMMA LEMMA LEMMA
\begin{lemma}\sl \label{lem:multiplication_V_endpoint_X_norm} If $ V\in  L^{n/2}(\R^n) $ then we have the H\"older-type inequality 
\[
\left |\int_{\R^n} V f g \right| \leq \|V\|_\lambda   \| f \|_{\X^*_\lambda} \| g \|_{\X^*_\lambda} 
\]
for all $f,g\in \X^*_\lambda$.  Consequently, we have that 
\[
\| V f \|_{\X_\lambda} \leq \|V\|_\lambda \| f \|_{\X_\lambda^*}.
\]
\end{lemma}
%%%%%%%%%%%%%%%%%%%%%%%%%%%%%% LEMMA LEMMA LEMMA

%%%%%%%%%%%%%%%%%%%%%%%%%%%%%% PROOF PROOF PROOF
\begin{proof}Since $1- \frac{2}{p_n} = \frac{2}{n} $ and $ 1 - \frac{2}{q_n} = \frac{2}{n+1} $, using \eqref{eq:holder} and Remark~\ref{rem:endpoint_norm_for_large_V} we get that
\[
\begin{split}
\Big| \int_{\R^n} V fg \Big| &\leq \Big| \int_{\R^n} \ind_{E_\lambda} V fg \Big| + \Big| \int_{\R^n} \ind_{\R^n \setminus E_\lambda} V fg \Big| 
\\
&\leq \| V \ind_{E_\lambda} \|_{L^{n/2}(\R^n)} \| f\|_{L^{p_n}(\R^n)} \| g\|_{L^{p_n}(\R^n)} 
+ \|V \ind_{\R^n \setminus E_\lambda} \|_{L^{\frac{n+1}{2}}(\R^n)} \| f \|_{L^{q_n}(\R^n)} \| g \|_{L^{q_n}(\R^n)} 
\\
&\leq \|V\|_\lambda\|f \|_{\X^\ast_\lambda} \| g \|_{\X^\ast_\lambda}
\end{split}
\]
as desired.
\end{proof}
%%%%%%%%%%%%%%%%%%%%%%%%%%%%%% PROOF PROOF PROOF

The spaces $X_\lambda$ and $X_\lambda^*$ are designed so that $P_\lambda\circ V$ regains a smallness property at the critical exponent $q=n/2$. This yields the promised refinement of \eqref{eq:KRS} in the present setting. We now turn to the verification of \eqref{eq:modRKS} for $X_\lambda^*$.

%%%%%%%%%%%%%%%%%%%%%%%%%%%%%% LEMMA LEMMA LEMMA
\begin{lemma}\sl \label{lem:modRKS} Let $n\geq 3$ and $\lambda>0$. For $X_\lambda$ as defined above, there holds
\[
\|P_\lambda f\|_{X_\lambda ^*} \lesssim \|f\|_{X_\lambda}.
\]
\end{lemma}
%%%%%%%%%%%%%%%%%%%%%%%%%%%%%% LEMMA LEMMA LEMMA

%%%%%%%%%%%%%%%%%%%%%%%%%%%%%% PROOF PROOF PROOF
\begin{proof} The conclusion of the lemma is equivalent to the four estimates below
\[
\begin{split}
\lambda^{\frac{1}{n+1}}\|P_\lambda f\|_{L^{q_n}(\R^n)}\lesssim  \lambda^{-\frac{1}{n+1}}\|f\|_{L^{q_n '}(\R^n)},&\qquad \|P_\lambda f\|_{L^{p_n}(\R^n)}\lesssim \|f\|_{L^{p_n '}(\R^n)},
\\
\lambda^{\frac{1}{n+1}}\|P_\lambda f\|_{L^{q_n}(\R^n)}\lesssim \|f\|_{L^{p_n '}(\R^n)},&
\qquad \|P_\lambda f\|_{L^{p_n}(\R^n)}\lesssim  \lambda^{-\frac{1}{n+1}}\|f\|_{L^{q_n '}(\R^n)}.
\end{split}
\]
The estimates in the first line of the display above are exactly the endpoint cases of \eqref{eq:KRS}. The two estimates in the second line are dual to each other so it will suffice to prove
\begin{equation}\label{eq:RKSsdual}
      \lambda^{\frac{1}{n+1}}\|P_\lambda f\|_{L^{p_n}(\R^n)}\lesssim \|f\|_{L^{q_n '}(\R^n)}.
\end{equation}

We turn to this task now. Begin by choosing a smooth bump function $\vp$ on $\R^n$ with $0\leq \vp\leq 1$, $\vp \equiv 1$ in $\{ x\in \R^n : \, |x|<2 \}$ and $\supp \vp \subset \{x\in\R^n: \, |x|< 4\}$. We fix $f\in \mathcal S(\R^n)$ and define the smooth frequency projections
\[
\wh{S_{<\lambda}f}(\xi) \coloneqq \vp ( \xi/\lambda ) \wh{f}(\xi),\qquad \wh{S_{>\lambda}f}(\xi) \coloneqq [ 1 - \vp ( \xi/\lambda ) ] \wh{f}(\xi),\qquad \xi\in\R^n.
\]
Let us write also $S_k f \coloneqq S_{<2^{k+1}}f-S_{<2^k}f$ for $k\in\Z$. We estimate
\begin{equation}\label{eq:resolvent_projections_split}
\|P_\lambda f \|_{L^{p_n}(\R^n)} \leq \| P_\lambda(S_{<\lambda} f) \|_{L^{p_n}(\R^n)} + \| P_\lambda(S_{>\lambda}f) \|_{L^{p_n}(\R^n)}.
\end{equation}
Since $P_\lambda(S_{<\lambda}f)$ is supported in frequency in $B(0,4\lambda)$ and
$p_n>q_n$, Bernstein's inequality gives
\begin{align*}
 \| P_\lambda(S_{<\lambda} f) \|_{L^{p_n}(\R^n)} &\lec (\lambda^n)^{\frac{1}{q_n} - \frac{1}{p_n}}  \| P_\lambda(S_{<\lambda} f) \|_{L^{q_n}(\R^n)} \lec \lambda^{\frac{1}{n+1}} \| P_\lambda f \|_{L^{q_n}(\R^n)}\lesssim \lambda^{- \frac{1}{n+1}} \| f \|_{L^{q_n'}(\R^n)}.
 \end{align*}

For the second summand in \eqref{eq:resolvent_projections_split} note that $|\xi|>2\lambda$ on the frequency support of $\wh{S_{>\lambda} f}$, which means that
$ \| \widehat{P_\lambda(S_k f)} \|_{L^2(\R^n)} = 2^{2k} \| \widehat{S_k f} \|_{L^2(\R^n)}$
whenever $2^k > 2 \lambda$. Then we use again that $q_n ' \leq 2 \leq p_n$ together with Bernstein's inequality, symbol estimates and again Bernstein's inequality, to get  
\[
\begin{split}
 \| P_\lambda (S_{>\lambda} f) \|_{L^{p_n}(\R^n)} &\lec \sum_{2^k > 2 \lambda} 2^{kn(\frac{1}{2} - \frac{1}{p_n})}      \| P_\lambda(S_k f) \|_{L^2(\R^n)} 
\simeq \sum_{ 2^k\gtrsim \lambda} 2^{-k} \| S_k f \|_{L^2(\R^n)} 
\\
&\lec \sum_{2^k \gtrsim \lambda} 2^{-k} 2^{kn(\frac{1}{q_n'} - \frac{1}{2})}  \| S_k f \|_{L^{q_n'}(\R^n)}
\lesssim \lambda^{-\frac{1}{n+1}} \| f \|_{L^{q_n'}(\R^n)}.
\end{split}
\]
This yields the desired estimate for the second summand in \eqref{eq:resolvent_projections_split}. Summing the estimates for the low and high frequencies of $P_\lambda f$ proves \eqref{eq:RKSsdual} and completes the proof of the lemma.
\end{proof}
%%%%%%%%%%%%%%%%%%%%%%%%%%%%%% PROOF PROOF PROOF

Combining the previous results and the discussion in this subsection it is now routine to fill in the details of the proof of the following corollary, where the perturbations $w^{\mathrm{cor}}$ are constructed for critical potentials $V\in L^1(\R^n)\cap L^{n/2}(\R^n)$.


%%%%%%%%%%%%%%%%%%%%%%%%%%%%%% COROLLARY COROLLARY COROLLARY
\begin{corollary}\label{cor:solutionsVcrit}\sl Let $n\geq 3$ and consider a time-independent potential $V\in L^1(\R^n)\cap L^{n/2}(\R^n)$.  Let $X_\lambda,X_\lambda ^*$ and $\|V\|_\lambda$ as defined above. There exists a constant $\lambda_V$ depending only on $V$ and the dimension such that for $\lambda > \lambda_V$, the functions
\[
w(x)\coloneqq e^{-i\lambda\omega\cdot x} +w^{\mathrm{cor}} (x)\eqqcolon w^{(0)}(x)+w^{\mathrm{cor}}(x),\qquad w^{\mathrm{cor}}(x)\coloneqq (\mathrm{Id}-P_\lambda\circ V)^{-1}[P_\lambda(Vw^{(0)})],
\]
with $\omega\in\mathbb S^{n-1} $ and $ x\in\R^n$,
are weak solutions of the time-independent Schr\"odinger equation 
\[
(\Delta+\lambda^2-V)w=0\quad\mathrm{in}\quad\R^n.
\]
Additionally, the following estimates are satisfied
\[
\left\|w^{\mathrm{cor}} \right\|_{X_\lambda ^*}\lesssim \|V\|_{X_\lambda}\leq \lambda^{-\frac{1}{n+1}} \|V\|_{L^{q_n '}(\R^n)}\leq  \lambda^{-\frac{1}{n+1}} \|V\|_{L^1(\R^n)\cap L^{n/2}(\R^n)}
\]
In particular, the function 
\[
\psi_V ^{\lambda,\omega}(t,x)\coloneqq e^{-i\lambda^2 t}w(x)= e^{-i\lambda^2 t}(w^{\mathrm{cor}}(x)+w^{(0)}(x)),\qquad (t,x)\in\R\times \R^n,
\]
is a weak solution of \eqref{sch_timeind}.
\end{corollary}
%%%%%%%%%%%%%%%%%%%%%%%%%%%%%% COROLLARY COROLLARY COROLLARY

%%%%%%%%%%%%%%%%%%%%%%%%%%%%%% PROOF PROOF PROOF
\begin{proof} As discussed in the comments preceding this lemma, we will invert the operator $(\mathrm{Id}-P_\lambda \circ V)$ on $\mathcal L(X_\lambda ^*)$. We combine the estimates of Lemma~\ref{lem:multiplication_V_endpoint_X_norm} with the one of Lemma~\ref{lem:modRKS} to estimate for each 
$F\in X_\lambda ^*$ 
\[
\|(P_\lambda \circ V )F\|_{X_\lambda ^*} \lesssim \|VF\|_{X_\lambda}\leq \|V\|_\lambda \|F\|_{X_\lambda ^*}.
\]
Now Remark~\ref{rem:endpoint_norm_for_large_V} shows that $\|P_\lambda \circ V \|_{\mathcal L(X_\lambda ^*)}<\frac{1}{2}$ if $\lambda$ is sufficiently large, say $\lambda>\lambda_V$ for some constant $\lambda_V$ depending on $V$. We can thus invert $\mathrm{Id}-P_\lambda\circ V$ in $\mathcal L(X_\lambda ^*)$ for all $\lambda>\lambda _V$ and $\|(\mathrm{Id}-P_\lambda\circ V)^{-1}\|_{\mathcal L(X_\lambda ^*)}<2$. Consequently, the following chain of estimates holds
\[
\|w^{\mathrm{cor}} \|_{X_\lambda ^*}\lesssim \|P_\lambda(Vw^{(0)})\|_{X_\lambda ^*}\lesssim \|V\|_{X_\lambda}.
\]
In order to complete the proof we just notice that $L^{q_n '}$ embeds continuously into $X_\lambda$ and the following norm estimate holds
\[
\| f\|_{X_\lambda} \leq \lambda^{-\frac{1}{n+1}}\|f\|_{L^{q_n '}}
\]
which yields the desired estimate when applied to $V$. Note that $1<q_n '=\frac{2(n+1)}{n+3}<n/2$ for $n\geq 3$ so the $L^{q_n '}$ norm of $V$ is controlled by $\|V\|_{L^1(\R^n)\cap L^{n/2}(\R^n)}$ and the proof is complete.
\end{proof}
%%%%%%%%%%%%%%%%%%%%%%%%%%%%%% PROOF PROOF PROOF


%%%%%%%%%%%%%%%%%%%%%%%%%%%%%% SECTION SECTION SECTION
\section{An orthogonality formula for stationary states}\label{sec:orthorelation} The goal of this section is to prove an orthogonality formula in the spirit of \eqref{eq:orthogonality_intro}, but one where the physical solutions $u_1,v_2$ are replaced by stationary state solutions such as those constructed in Section~\ref{sec:timeindependent_solutions} and in particular having the form \eqref{eq:steadysol}. We first need to state and prove a series of integration by parts identities which will be used in extending the orthogonality relation \eqref{eq:orthogonality_intro}. 


%%%%%%%%%%%%%%%%%%%%%%%%%%%%%% SECTION SECTION SECTION
\subsection{Integration by parts identities} We record several integration by parts formulas that will be used to derive the orthogonality relation for stationary states in Proposition~\ref{pr:int_by_parts_time-harmonic_solutions}.

We first state a special case of \cite{caro2025initialtofinalstateinverseproblemunbounded}*{Proposition 5.2}, adapted to our setting. The proof is omitted, since it follows the same argument; see \cite{caro2025initialtofinalstateinverseproblemunbounded}*{Appendix A} for details of the proof.

%%%%%%%%%%%%%%%%%%%%%%%%%%%%%% PROPOSITION PROPOSITION PROPOSITION
\begin{proposition}\label{prp:integration_by_parts_unbounded_no-endpoint}\sl Let $n\geq 2$ and $(r,p)$ be a Strichartz pair. Assume
\[
u,v \in C\left ([0,T];L^2(\R^n)\right) \cap L^r\left((0,T);L^p(\R^n)\right)
\]
be such that
\[
(i\partial_t+\Delta)u, (i\partial_t+\Delta)v \in L^{r'}\big((0,T);L^{p'}(\R^n)\big).
\]
Then
\[
\int_{\Sigma} \left[(i\partial_t+\Delta)u \ovl v - u \ovl{(i\partial_t+\Delta)v}\right] = i \int_{\R^n} \left[u(T,\cd)\ovl{v(T,\cd)}-u(0,\cd)\ovl{v(0,\cd)}\right].
\]
\end{proposition}
%%%%%%%%%%%%%%%%%%%%%%%%%%%%%% PROPOSITION PROPOSITION PROPOSITION

We now prove an extension of Proposition~\ref{prp:integration_by_parts_unbounded_no-endpoint} in the special case $u(0,\cd)=u(T,\cd)=0$, where the integration by parts identity remains valid even when the second function is only assumed to lie in $C([0,T];L^{p}(\R^{n}))$.

%%%%%%%%%%%%%%%%%%%%%%%%%%%%%% PROPOSITION PROPOSITION PROPOSITION
\begin{proposition}\label{cor:integration_by_parts_unbounded_no-endopint2}\sl
Let $n\geq 2$ and $(r,p)$ be a Strichartz pair. Let
\[
v \in C\left([0,T];L^p(\R^n)\right)
\]
and
\[
u \in C\left([0,T];L^2(\R^n)\right) \cap L^r\left((0,T);L^p(\R^n)\right)\quad\text{with}\quad u(0,\cd)=u(T,\cd)=0,
\]
be such that
\[
\left(i\partial_t+\Delta\right)u, \left(i\partial_t+\Delta\right)v \in L^{r'}\big((0,T);L^{p'}(\R^n)\big).
\]
Then
\[
\int_\Sigma \left(i\partial_t+\Delta\right)u \ovl v = \int_\Sigma u \ovl{\left(i\partial_t+\Delta\right)v}.
\]
\end{proposition}
%%%%%%%%%%%%%%%%%%%%%%%%%%%%%% PROPOSITION PROPOSITION PROPOSITION


%%%%%%%%%%%%%%%%%%%%%%%%%%%%%% PROOF PROOF PROOF
\begin{proof} If we additionally assume that $v\in C([0,T];L^2(\R^n))$, the conclusion follows directly from Proposition~\ref{prp:integration_by_parts_unbounded_no-endpoint}, since the boundary terms vanish due to the conditions $u(0,\cd)=u(T,\cd)=0$. The main part of the proof is therefore devoted to removing the assumption $v\in C([0,T];L^2(\R^n))$ and reducing the case $v\in C([0,T];L^p(\R^n))$ to the $L^2$-case via approximation.

Let $\chi$ be a smooth bump function on $\R^n$ such that $0\leq \chi \leq 1$, $\chi\equiv 1$ on $\{|x|\leq 1\}$, and $\supp \chi \subset \{|x|\leq 2\}$. Let also $\phi\in \mathcal S(\R^n)$ be a nonnegative bump function with $\supp \phi \subset B(0,1)$, satisfying $\int_{\R^n}\phi=1$. For $\varepsilon>0$, let $R=R(\varepsilon)>0$ be a function such that
\[
\lim_{\varepsilon\to0} R(\varepsilon)=+\infty,
\]
which will be chosen in the course of the proof.

We define the rescaled bump functions $\chi^R(x)\coloneqq \chi(x/R)$ and $\phi_\varepsilon(x)\coloneqq \varepsilon^{-n}\phi(x/\varepsilon)$ and for general $v \in C([0,T];L^p(\R^n))$, we set
\[
v_{R,\varepsilon} \coloneqq \chi^R (\phi_\varepsilon * v).
\]
Here the convolution $\phi_\varepsilon * v$ is taken only in the space variables.  Since $v\in C([0,T];L^p(\R^n))$ with $p\geq 2$ and $\chi^R$ has compact support, it follows that
\[
v_{R,\varepsilon}\in C\left([0,T];L^2(\R^n)\right) \cap L^r\left((0,T);L^p(\R^n)\right).
\]
Now an application of the Leibniz rule in the sense of distributions yields
\begin{equation}\label{eq:Leibniz_distributional}
\left(i \d_t + \Delta \right)v_{R, \ve} = \chi^R \left(i\d_t + \Delta\right)(v*\phi_\ve) +  ( v*\phi_\ve ) \Delta(\chi^R)  + 2 ( \nabla \chi^R ) \cdot \nabla (v*\phi_\ve),
\end{equation}
in $\mathcal{D}' (\Sigma)$. We record for later use the easy identities
\begin{equation}\label{eq:scalederiv}
\begin{split}
& \Delta(\chi^R) = R^{-2} (\Delta \chi)^R \ind_{\{R\leq |x| \leq 2R\}}, \qquad \nabla (\chi^R) = R^{-1} (\nabla \chi)^R \ind_{\{R\leq |x| \leq 2R\}}, 
\\
&  \nabla (\phi_\ve) = \ve^{-1} (\nabla \phi)_{\ve}. 
\end{split}
\end{equation}

We estimate each term in \eqref{eq:Leibniz_distributional}. First, we have
\[
\chi^R (i\d_t + \Delta)(v*\phi_\ve) = \chi^R \phi_\ve * (i\d_t + \Delta)v \in L^{r'}\big((0, T); L^{p'}(\R^n)\big).
\]
Moreover, the function
\[
( v*\phi_\ve ) \Delta(\chi^R) = R^{-2} (\Delta \chi)^R (v*\phi_\ve)
\]
has support contained in $\{x\in\R^n:\, R\leq |x|\leq 2R\}$ by the definition of $\chi$. Since $p\geq 2$ we have that $p'\leq 2 \leq p$, and hence $(v*\phi_\ve) \Delta (\chi^R)\in C([0, T]; L^{p'}(\R^n))$.  Similarly, we also infer that 
\[
( \nabla \chi^R ) \nabla (v*\phi_\ve) = R^{-1} \ve^{-1} (\nabla \chi)^R (v* (\nabla \phi)_\ve) \in C\big ([0, T]; L^{p'}(\R^n)\big). 
\]
Therefore, $(i \d_t + \Delta )v_{R, \ve} \in L^{r'}((0, T); L^{p'}(\R^n))$ and we can use the integration by parts formula of Proposition~\ref{prp:integration_by_parts_unbounded_no-endpoint} with the right hand side equal to zero to get 
\[
\int_\Sigma [(i\d_t + \Delta)u] \ovl{v_{R,\ve}} = \int_\Sigma \ovl{(i\d_t + \Delta) v_{R, \ve}}  u.
\]
Notice that, in order to complete the proof it will be enough to show that 
\begin{equation}\label{eq: Cor_int_by_parts_approx1}
\lim_{\ve \to 0} \int_\Sigma \left[(i\d_t + \Delta) u\right] \ovl{(v_{R, \ve} - v)} = 0, 
\end{equation}
and
\begin{equation}\label{eq: Cor_int_by_parts_approx2}
\lim_{\ve \to 0} \int_\Sigma u  \ovl{\left(i\d_t + \Delta\right)(v_{R, \ve} - v)} = 0.
\end{equation}

We begin with the proof of~\eqref{eq: Cor_int_by_parts_approx1}. We have that 
\begin{align*}
\Big| \int_\Sigma  [(i\d_t + \Delta) u] \ovl{(v_{R, \ve} - v)} \Big| &\leq 
\| (i\d_t + \Delta)u \|_{L^{r'}((0, T); L^{p'}(\R^n))} \| v_{R, \ve} - v \|_{L^r((0, T); L^p(\R^n))} \\
&\leq \| (i\d_t + \Delta)u \|_{L^{r'}((0, T); L^{p'}(\R^n))} \| (\chi^R - 1) v\|_{L^r((0, T); L^p(\R^n))} \\
&+ \| (i\d_t + \Delta)u \|_{L^{r'}((0, T); L^{p'}(\R^n))} \| \phi_\ve * v - v \|_{L^r((0, T); L^p(\R^n))}
\end{align*}
which implies~\eqref{eq: Cor_int_by_parts_approx1} as $\ve \to 0$ by dominated convergence. 

We now turn to the proof of the more challenging limit~\eqref{eq: Cor_int_by_parts_approx2}. We have that 
\begin{equation}\label{eq: Cor_int_by_parts_approx2_split} 
\begin{split}
\left| \int_\Sigma \big[ \ovl{ (i\d_t + \Delta)(v_{R,\ve} -v) } \big] u \right| &\leq
\left| \int_\Sigma \big[ \ovl{ \chi^R [(i\d_t+\Delta)v]*\phi_\ve - (i\d_t+\Delta)v }\big] u \right| 
\\
&\qquad +\left| \int_\Sigma \overline{\Delta(\chi^R) (v*\phi_\ve)} u \right| 
+ 2 \left| \int_\Sigma \overline{\nabla (\chi^R) \cdot \nabla (v*\phi_\ve)} u \right|
\\
&\eqqcolon \mathrm{I}(\ve,R)+\mathrm{II}(\ve,R)+\mathrm{III}(\ve,R).
\end{split}  
\end{equation}

We begin by estimating $\mathrm I(\ve,R)$. We have
\[
\begin{split}
\mathrm{I}(\ve,R) &\leq  \int_\Sigma \left| \left[\ovl{\phi_\ve*\left(i\d_t+\Delta\right)v - \left(i\d_t+\Delta\right)v } \right] \chi^R \right| |u|  + \int_\Sigma \left|\chi^R - 1\right| \left|\left(i\d_t +\Delta\right)v \right| |u| 
\\
&\leq \left\| \phi_\ve*\left(i\d_t+\Delta\right)v - \left(i\d_t+\Delta\right)v \right\|_{L^{r'}((0, T); L^{p'}(\R^n))} \| u \|_{L^r((0, T); L^p(\R^n))} 
\\
& \qquad + \int_\Sigma \left|\chi^R - 1\right| \left|\left(i\d_t +\Delta\right)v \right| |u|.
\end{split}
\]
In the right-hand side of the display above, the first summand tends to $0$ as $\varepsilon\to0$ by the convergence of the approximate identity $\phi_\varepsilon*(i\d_t+\Delta)v$ in $L^{p'}(\R^n)$, while the second summand tends to $0$ as $R\to\infty$ by dominated convergence. Consequently, $\mathrm I(\ve,R(\ve))\to 0$ as $\varepsilon\to0$ with any choice of $R=R(\varepsilon)\to\infty$ as $\ve \to 0$.

We estimate the worst term $\mathrm{III}(\varepsilon,R)$. By the assumptions on $\chi,\phi$ and \eqref{eq:scalederiv}, we have
\[
\left|\nabla(\chi^R)\cdot\nabla(v*\phi_\varepsilon)\right|= \ve^{-1} \left|\nabla(\chi^R)\cdot (v*(\nabla\phi)_\varepsilon)\right| \lesssim R^{-1}\varepsilon^{-1} \ind_{\{R\le |x|\le 2R\}}  \left(|v|*|(\nabla\phi)_\varepsilon|\right).
\]
Since $\supp (\nabla\phi)_\varepsilon \subset B(0,\varepsilon)$, it follows that, provided
$\varepsilon\ll R$,
\[
\ind_{\{R\le |x|\le 2R\}}\left(|v|*|(\nabla\phi)_\varepsilon|\right) \leq \left(\ind_{\{|x|\simeq R\}}|v|\right)*|(\nabla\phi)_\varepsilon|.
\]
Therefore,
\[
\begin{split}
\mathrm{III}(\varepsilon,R) &\lesssim R^{-1}\varepsilon^{-1}\int_\Sigma \left[\left(\ind_{\{|x|\simeq R\}}|v|\right)*|(\nabla\phi)_\varepsilon|\right](t,x) |u(t,x)|\,\dd x\,\dd t
\\
&\le R^{-1}\varepsilon^{-1}  \left\|\left(\ind_{\{|x|\simeq R\}}|v|\right)*|(\nabla\phi)_\varepsilon|\right\|_{L^1((0,T);L^2(\R^n))} \|u\|_{C([0,T];L^2(\R^n))}
\\
&\le R^{-1}\varepsilon^{-1} \|v\|_{L^1((0,T);L^2(\{|x|\simeq R\}))}\|u\|_{C([0,T];L^2(\R^n))} 
\\
&\lesssim R^{-1}\varepsilon^{-1}  R^{n(\frac12-\frac1p)} \|v\|_{L^1((0,T);L^p(\{|x|\simeq R\}))}\|u\|_{C([0,T];L^2(\R^n))} .
\end{split}
\]
Here, we used Cauchy--Schwarz in $x$ and H\"older in $t$ to pass to the second line, Young's inequality in $x$ in passing to the third line, together with $\|(\nabla\phi)_\varepsilon\|_{L^1(\R^n)}\lesssim 1$, and finally H\"older on the annulus $\{|x|\simeq R\}$ to pass from $L^2(\{|x|\simeq R\})$ to $L^p(\{|x|\simeq R\})$ with $p\geq2$. Using the calculation
\[
-1+n\left(\frac12-\frac1p\right)= n\left(\frac1{p_n}-\frac1p\right), \qquad  \frac1{p_n}=\frac12-\frac1n,
\]
and writing $R=R(\varepsilon)$, we have proved
\[
\mathrm{III}(\varepsilon,R(\varepsilon)) \lesssim \varepsilon^{-1}R(\varepsilon)^{-n(\frac1p-\frac1{p_n})} \|v\|_{L^1((0,T);L^p(\{|x|\simeq R(\varepsilon)\}))} \|u\|_{C([0,T];L^2(\R^n))}.
\]

If $p<p_n$, then $n(\frac1p-\frac1{p_n})>0$, and choosing
\[
R(\varepsilon)\coloneqq \varepsilon^{-\frac{2}{n(\frac1p-\frac1{p_n})}}
\]
gives $\mathrm{III}(\varepsilon,R(\varepsilon))\lesssim \varepsilon$, hence $\mathrm{III}(\varepsilon,R(\varepsilon))\to0$ as $\varepsilon\to0$.

In the endpoint case $p=p_n$ we have to be more careful since we get that
\[
\mathrm{III}(\ve,R(\ve))
\leq \frac{1}{\ve}  \|u\|_{C([0, T]; L^2(\R^n))}  \|v\|_{L^1\left((0, T); L^{ p_n}(|x|\simeq R(\ve))\right)} .
\]
Notice that
\[
v\ind_{\{|x|\simeq R\}}\ind _{(0,T)} \in C\left([0,T];L^{p_n}(\R^n)\right) \subset L^1\left((0,T);L^{p_n}(\R^n)\right),
\]
and that
\[
v(t,x)\ind_{\{|x|\simeq R\}}\ind_{(0,T)}(t)\to0 \quad\text{as }R\to\infty
\]
for every $(t,x)\in\Sigma$. By dominated convergence,
\[
\lim_{R\to\infty} \|v\|_{L^1((0,T);L^{p_n}(\{|x|\simeq R\}))}=0.
\]
Hence, for every $\ve>0$ there exists $R(\ve)>0$ which can be chosen to also satisfy $R(\ve)>\ve^{-1}$ such that $\|  v \|_{L^1((0,T); L^{p_n}(\{|x|\simeq R\}))}<\ve ^2$ for all $R\geq R(\ve)$. This shows that in the endpoint case $p=p_n$ one can choose $R(\ve)$, with $\lim_{\ve\to0 }R(\ve)=+\infty$, so that
\[
\mathrm{III}(\ve,R(\ve)) \lesssim \frac{1}{\ve}  \|u\|_{C([0, T]; L^2(\R^n))}  \|v\|_{L^1((0, T); L^{p_n}(|x|\simeq R(\ve)))} \to 0 \quad \text{as}\quad \ve\to 0.
\]

The term $\mathrm{II}(\varepsilon,R(\varepsilon))$ is estimated similarly. In particular, in the endpoint case $p=p_n$ we obtain
\[
\mathrm{II}(\ve,R(\ve)) \lesssim  \frac{1}{R(\ve)} \|u\|_{C([0, T]; L^2(\R^n))} \|v\|_{L^1((0, T); L^{p_n}(\R^n))} \to 0 \quad \textup{as} \quad \ve\to 0. 
\]
Combining the estimates we see that there is a choice $R(\ve)\to \infty$ when $\ve\to 0$ such that 
\[
\lim_{\ve\to 0} \left[\mathrm{I}(\ve,R(\ve))+\mathrm{II}(\ve,R(\ve))+\mathrm{III}(\ve,R(\ve))\right]=0
\]
so the right hand side of~\eqref{eq: Cor_int_by_parts_approx2_split} tends to 0 as $\ve\to 0$. This  shows~\eqref{eq: Cor_int_by_parts_approx2} and completes the proof of  Proposition~\ref{cor:integration_by_parts_unbounded_no-endopint2}. 
\end{proof}
%%%%%%%%%%%%%%%%%%%%%%%%%%%%%% PROOF PROOF PROOF

With this in hand we proceed with the following integration by parts formula for complex exponential solutions. 

%%%%%%%%%%%%%%%%%%%%%%%%%%%%%% LEMMA LEMMA LEMMA
\begin{lemma}\label{lem:int_by_parts_complex-exp_Lp}\sl Let $(r,p)$ be a Strichartz pair with $p\ge q_n$, and assume in addition that $p\le p_n$ if $n\ge3$, or $p<\infty$ if $n=2$. Let
\[
u\in C\left([0,T];L^2(\R^n)\right )\cap L^r\left((0,T);L^p(\R^n)\right), \qquad u(0,\cd)=u(T,\cd)=0,
\]
and suppose that
\[
(i\partial_t+\Delta)u\in L^1(\Sigma)\cap L^{r'}\big((0,T);L^{p'}(\R^n)\big).
\]
Then
\[
\int_\Sigma e^{i(\lambda^2 t+\lambda\omega\cdot x)}  (i\partial_t+\Delta)u = 0,
\]
for every $\lambda>0$ and every $\omega\in\Sph^{n-1}$.
\end{lemma}
%%%%%%%%%%%%%%%%%%%%%%%%%%%%%% LEMMA LEMMA LEMMA

%%%%%%%%%%%%%%%%%%%%%%%%%%%%%% PROOF PROOF PROOF
\begin{proof} We first construct an approximation of the function $\psi(t,x)\coloneqq e^{  -i \lambda^2 t - i\lambda \omega \cdot x }$. Our favorite way of doing this is the following. Consider $\chi \in \mathcal S (\R^n)$ such that $0 \leq \chi (x) \leq 1$ for all $x \in \R^n$ and $\supp \chi \subset \{ x \in \R^n : \, |x| < 1/2 \}$. We will also choose $\chi$ so that 
\[
\chi(0)=\frac{1}{(2\pi)^{n/2}}\int_{\R^n} \widehat{\chi}(\eta)\,\dd \eta=1.
\]
For $\varepsilon > 0$, define
\[ 
f^\ve (x) \coloneqq \chi(\ve x) e^{-i \lambda \omega \cdot x}\eqqcolon \chi^\ve (x) e^{-i \lambda \omega \cdot x}, \qquad  x \in \R^n. 
\]
Clearly $f^\ve \in \mathcal S(\R^n)$ for each $\ve>0$. Now let $\psi_\ve: \R\times \R^n\to \C$ be defined as 
\[
\psi_\ve (t,x)\coloneqq e^{it\Delta}(f^\ve)(x)=\frac{1}{(2\pi)^{n/2}}\int_{\R^n} e^{-i|\eta|^2t + i \eta\cdot x} \widehat{f^\ve}(\eta)\, \dd \eta,\qquad (t,x)\in \R\times \R^n.
\]
By construction we have that $(i\partial_t+\Delta)\psi_\ve=0$ in $\R\times \R^n$ and $\psi_\ve \in C^\infty(\R;\mathcal S(\R^n))$. In particular $\psi_\ve\in C([0,T];L^2(\R^n))\cap L^{r}((0,T);L^p(\R^n))$ for any Strichartz pair $(r,p)$. 

Now, consider a function $u$ as in the assumption and apply Proposition~\ref{cor:integration_by_parts_unbounded_no-endopint2} to the functions $u,\psi_\ve$ as above to get
\[
 \int_{\Sigma}  \left(i\partial_t+\Delta\right)u  \ovl{\psi_\ve} = \int_{\Sigma} u  \ovl{\left(i\partial_t+\Delta\right)\psi_\ve}  = 0.     
\]
Noting that $\widehat {f^\ve}(\eta)=\widehat{\chi^\ve}(\eta+\lambda \omega)=\ve^{-n}\widehat{\chi}((\eta+\lambda\omega)/\ve)$, we write for $(t,x)\in \Sigma$
\[
\begin{split}
\psi_\ve (t,x)-\psi(t,x)& = \frac{1}{(2\pi)^{n/2}}\int_{\R^n} \widehat{f^\ve}(\eta) e^{-i|\eta|^2t + i\eta\cdot x}\, \dd \eta - e^{-i\lambda^2 t - i\lambda\omega\cdot x}
\\
&= \frac{1}{(2\pi)^{n/2}}\int_{\R^n}\widehat{\chi^\ve}(\eta+\lambda\omega) e^{-i|\eta|^2t+i\eta\cdot x}\, \dd \eta - e^{-i\lambda^2 t -i\lambda\omega\cdot x}
\\
&= \frac{1}{(2\pi)^{n/2}}\int_{\R^n}\frac{1}{\ve^n}\widehat{\chi}\left(\frac{\xi}{\ve}\right) e^{-i|\xi-\lambda\omega|^2t +i(\xi-\lambda\omega)\cdot x}\, \dd \xi - e^{-i\lambda^2 t - i\lambda\omega\cdot x}
\\
&=\frac{e^{-i\lambda^2 t - i\lambda\omega\cdot x}}{(2\pi)^{n/2}}
\int_{\R^n} \widehat{\chi }(\eta ) \left[e^{-i\left(\ve^2|\eta|^2-2\ve \lambda \omega\cdot \eta\right)t + i\ve\eta\cdot x} - 1\right]\, \dd \eta ,
\end{split}
\]
by our assumption $\chi(0)=1$. Since $\chi\in\mathcal S(\R^n)$, we have $\psi_\ve(t,x)\to \psi(t,x)$ as $\ve \to 0$ for every $(t,x)\in \R\times\R^n$ by dominated convergence. The conclusion now follows by yet another application of dominated convergence since
\[
\left| \ovl{(\psi_\ve-\psi)}{\left(i\partial_t+\Delta\right)u}\right|\leq \left(1+(2\pi)^{-n/2}\| \widehat{\chi} \|_{L^1(\R^n)}\right)\left|\left(i\partial_t+\Delta\right)u\right|\in L^1(\Sigma)
\]
since by our assumption $(i\partial_t+\Delta)u\in L^1(\Sigma)$.
\end{proof}
%%%%%%%%%%%%%%%%%%%%%%%%%%%%%% PROOF PROOF PROOF


%%%%%%%%%%%%%%%%%%%%%%%%%%%%%% SECTION SECTION SECTION
\subsection{Orthogonality formulas} We now relate the initial-to-final-state maps to the difference of the corresponding potentials via solutions of \eqref{eq:Schrodinger}. When $\mathcal U_T^1=\mathcal U_T^2$, this yields the basic orthogonality relation \eqref{eq:orthogonality_intro} for physical solutions.

%%%%%%%%%%%%%%%%%%%%%%%%%%%%%% PROPOSITION PROPOSITION PROPOSITION
\begin{proposition}\label{prp:CR_prop4.1_unbounded}\sl Let $V_1, V_2 \in L^q(\R^n)$ for some $q\in(1,\infty]$ if $n=2$, or for some $q\in[n/2,\infty]$ if $n\geq 3$, and let $(r, p)$ be a Strichartz pair with $\frac{1}{q}+\frac{2}{p}=1$.  For every fixed $T>0$ and every $f, g \in L^2(\R^n)$ we have that 
\[
i\int_{\R^n} (\cal{U}^1_T - \cal{U}^2_T) f\ovl{g} = \int_{\Sigma} (V_1 - V_2) u_1 \ovl{v_2},
\]
where $u_1 \in C([0, T]; L^2 (\R^n)) \cap L^r((0, T); L^p(\R^n))$ is a solution of \eqref{eq:Schrodinger} with potential $V_1$ and initial data $f$, while $v_2 \in C([0, T]; L^2 (\R^n)) \cap L^r((0, T); L^p(\R^n))$ is the physical solution of the following final-value problem 
\begin{equation}\label{eq:FVP}
\begin{cases}
i\d_t v_2 = -\Delta v_2 + \ovl{V_2} v_2 \quad &\mathrm{in} \quad \Sigma,
\vspace{.6em}
\\
v_2(T, \cd) = g  &\mathrm{in} \quad \R^n.
\end{cases}
\end{equation}
\end{proposition}
%%%%%%%%%%%%%%%%%%%%%%%%%%%%%% PROPOSITION PROPOSITION PROPOSITION

We omit the proof of the proposition above since it follows exactly the same lines as~\cite{caro2025initialtofinalstateinverseproblemunbounded}*{Proposition 5.1}, by using the integration by parts formula of Proposition \ref{prp:integration_by_parts_unbounded_no-endpoint} in our case. 

We record the following lemma, which converts an orthogonality condition on the source term F against physical final-state solutions into a vanishing final condition for the corresponding inhomogeneous solution.
%%%%%%%%%%%%%%%%%%%%%%%%%%%%%% LEMMA LEMMA LEMMA
\begin{lemma}\label{lem:CR_lemma4.4_Lp}\sl Consider $V\in L^q(\R^n)$ for some $q\in(1,\infty]$ if $n=2$, or for some $q\in[n/2,\infty]$ if $n\geq 3$, let $(r, p)$ be a Strichartz pair with $\frac{1}{q}+\frac{2}{p}=1$, and $F\in L^{r'}((0,T); L^{p'}(\R^n))$ with 
\[
\int_\Sigma F \ovl{v} = 0
\]
for every $v\in C([0, T]; L^2(\R^n)) \cap L^r((0, T); L^p(\R^n))$  which is the solution of the final-value problem 
\[
\begin{cases}
i\d_t v = - \Delta v + \ovl{V} v\quad &\textrm{in}\quad \Sigma,
\vspace{.6em}
\\
v(T, \cd)= g  \quad &\textrm{in}\quad \R^n.
\end{cases}
\]
with arbitrary $g\in L^2(\R^n)$. Then, the solution $u\in C([0, T]; L^2(\R^n)) \cap L^r((0, T); L^p(\R^n))$ of the problem
\[
\begin{cases}
\left(i\partial_t +\Delta\right)u-Vu=F\quad &\textrm{in}\quad \Sigma,
\vspace{.6em}
\\
u(0,\cd)=0, \quad &\textrm{in}\quad \R^n.
\end{cases}
\]
satisfies $u(T, \cd) = 0. $
\end{lemma}
%%%%%%%%%%%%%%%%%%%%%%%%%%%%%% LEMMA LEMMA LEMMA

The proof is in the same spirit as~\cite{caro2025initialtofinalstateinverseproblemunbounded}*{Lemma 5.3}, but we include it also here for reader's convenience since it is rather short. 


%%%%%%%%%%%%%%%%%%%%%%%%%%%%%% PROOF PROOF PROOF
\begin{proof} Let $g\in L^2(\R^n)$ and $v$ be the corresponding solution of the final-value problem in the statement of the lemma. Under these assumptions, Proposition~\ref{prp:integration_by_parts_unbounded_no-endpoint} applies to $v$ yielding
\begin{align*}
i \int_{\R^n} u(T, \cd) \ovl{g} &= \int_\Sigma \left[ \left(i \d_t +\Delta\right)u \ovl{v} - u \ovl{(i\d_t + \Delta)v} \right]
= \int_\Sigma \left[ \left(i \d_t +\Delta - V\right)u\ovl{v} - u \ovl{\left(i \d_t + \Delta - \ovl{V}\right)v}  \right] 
\\
&= \int_\Sigma F \ovl{v}  = 0.
\end{align*}
Since this holds for every $g\in L^2(\R^n)$, we  conclude that $u(T, \cd) = 0$. 
\end{proof}
%%%%%%%%%%%%%%%%%%%%%%%%%%%%%% PROOF PROOF PROOF

The following lemma is a symmetric version of Lemma~\ref{lem:CR_lemma4.4_Lp}; we omit the completely analogous proof.

%%%%%%%%%%%%%%%%%%%%%%%%%%%%%% LEMMA LEMMA LEMMA
\begin{lemma}\label{lem:CR_lemma4.4_Lp_symmetric}\sl Consider $V\in  L^q(\R^n)$ for some $q\in(1,\infty]$ if $n=2$, or for some $q\in[n/2,\infty]$ if $n\geq 3$, let $(r, p)$ be a Strichartz pair with $\frac{1}{q}+\frac{2}{p}=1$, and $G\in L^{r'}((0,T); L^{p'}(\R^n))$ such that 
\[
\int_\Sigma \ovl{G} u = 0
\]
for every $u \in C([0, T]; L^2(\R^n)) \cap L^r((0, T); L^p(\R^n)) $ which is the solution of the initial-value problem \eqref{eq:Schrodinger} for arbitrary $f\in L^2(\R^n)$. Then, the solution $v \in C([0, T]; L^2(\R^n)) \cap L^r((0, T); L^p(\R^n))$ of the problem 
\[
\begin{cases}
\left(i\d_t + \Delta - \ovl{V}\right) v = G \quad &\textrm{in}\quad \Sigma, 
\vspace{.6em}\\
v(T, \cd) = 0  \quad &\textrm{in}\quad \R^n.
\end{cases} 
\]
satisfies that $v(0, \cd)=0$.
\end{lemma}
%%%%%%%%%%%%%%%%%%%%%%%%%%%%%% LEMMA LEMMA LEMMA

We are now in a position to prove the main orthogonality relation of this paper, in which the physical solutions appearing in Proposition~\ref{prp:CR_prop4.1_unbounded} are replaced by the stationary state solutions constructed in Section~\ref{sec:timeindependent_solutions}.

%%%%%%%%%%%%%%%%%%%%%%%%%%%%%% PROPOSITION PROPOSITION PROPOSITION
\begin{proposition}\label{pr:int_by_parts_time-harmonic_solutions}\sl Let $V_1,V_2\in L^1(\R^n)\cap L^q(\R^n)$, where $q\in(1,(n+1)/2]$ if $n=2$, or $q\in[n/2,(n+1)/2]$ if $n\geq 3$, and define $p$ by 
\[
  \frac{1}{q}+\frac{2}{p}=1.
\]
Let $\mathcal U_T^1,\mathcal U_T^2$ denote the initial-to-final-state maps corresponding to $V_1,V_2$, respectively.

Assume that $\mathcal U_T^1=\mathcal U_T^2$. Then, for every $\lambda>0$, every $\omega_1,\omega_2\in\Sph^{n-1}$, and every pair of functions
\[
\psi_j(t,x)=e^{-i\lambda^2 t}\left(w_j^{(0)}(x)+w_j^{\mathrm{cor}}(x)\right),
\qquad j\in\{1,2\},
\]
with
\[
w_j^{(0)}(x)=e^{-i\lambda\omega_j\cdot x},  \qquad w_j^{\mathrm{cor}}\in L^p(\R^n),
\]
satisfying
\[
\left(i\partial_t+\Delta-V_1\right)\psi_1=0, \qquad \left(i\partial_t+\Delta-\overline{V_2}\right)\psi_2=0,
\]
the following orthogonality relation holds:
\begin{equation}\label{eq:int_by_parts_time-harmonic} \int_\Sigma (V_1-V_2) \psi_1 \overline{\psi_2}=0.
\end{equation}

In particular, \eqref{eq:int_by_parts_time-harmonic} holds for the stationary-state solutions constructed in Corollaries~\ref{cor:solutions} and~\ref{cor:solutionsVcrit}.
\end{proposition}
%%%%%%%%%%%%%%%%%%%%%%%%%%%%%% PROPOSITION PROPOSITION PROPOSITION




%%%%%%%%%%%%%%%%%%%%%%%%%%%%%% PROOF PROOF PROOF
\begin{proof}  We argue by combining the orthogonality relation for physical solutions with suitable auxiliary boundary value problems in order to obtain \eqref{eq:int_by_parts_time-harmonic}. This is implemented in two steps. 

\vspace{1em}
\noindent\textbf{Step 1:} In this first step, we prove that
\[
\cal{U}^1_T=\cal{U}^2_T \Rightarrow \int_\Sigma (V_1 - V_2) \Psi \ovl{\psi_2} = 0 ,
\]
for any $\Psi \in C([0, T]; L^2(\R^n)) \cap L^r((0, T); L^p(\R^n))$ satisfying $(i \d_t + \Delta - V_1)\Psi = 0$ in $\Sigma$. 

To that end, for any such solution \(\Psi\) as above, consider the problem
\begin{equation}\label{eq:initial_BVP_F}
\begin{cases}
\left(i \d_t + \Delta - V_2\right )\Phi = (V_1 - V_2)\Psi \eqqcolon F  \quad &\textrm{in}\quad \Sigma, 
\vspace{.6em}\\
\Phi(0, \cd) = 0  \quad &\textrm{in}\quad \R^n.
\end{cases}
\end{equation}
Using that $V_1-V_2 \in L^q(\R^n)$ and \eqref{eq:holder1} we have $F\in L^{r'}((0,T);L^{p'}(\R^n))$, and hence, by the discussion in \S\ref{subsec: Stri_WP}, there exists a unique solution \(\Phi\) of \eqref{eq:initial_BVP_F} satisfying
\[
\Phi\in C\left([0,T];L^2(\R^n)\right)\cap L^r\left((0,T);L^p(\R^n)\right).
\] 
Now, since $\mathcal U_T^1=\mathcal U_T^2$, Proposition~\ref{prp:CR_prop4.1_unbounded}
implies that
\[
\int_\Sigma (V_1-V_2)\Psi \overline{v_2}=0
\]
for every physical final-state solution $v_2$ associated with $\ovl{V_2}$. In particular, this holds for the physical solution $v_2$ of the final-value problem with potential $\ovl{V_2}$ and final data $g=\Phi(T,\cd)$. Applying Lemma~\ref{lem:CR_lemma4.4_Lp} with $u=\Phi$ and $V=\ovl{V_2}$, we conclude that
$\Phi(T,\cd)=0$.

We now plan to apply Lemma~\ref{lem:int_by_parts_complex-exp_Lp} with $\Phi$ in the place of $u$. For this, we need to verify that
\[
\left(i\partial_t+\Delta\right)\Phi \in L^1(\Sigma)\cap L^{r'}\left((0,T);L^{p'}(\R^n)\right) \quad\text{and}\quad \Phi(0,\cd)=\Phi(T,\cd)=0.
\]
The boundary condition $\Phi(0,\cd)=0$ holds by construction, while we already proved that $\Phi(T,\cd)=0$. Moreover, from \eqref{eq:initial_BVP_F} we have
\[
\left(i\partial_t+\Delta\right)\Phi = V_2\Phi + F,
\]
so it suffices to verify that $V_2\Phi$ and $F$ belong to $L^1(\Sigma)\cap L^{r'}((0,T);L^{p'}(\R^n))$.

First, we check the $L^{r'}((0,T);L^{p'}(\R^n))$-bounds. Since $V_1-V_2,\,V_2\in L^q(\R^n)$ and $\Psi,\Phi\in L^{r}((0,T);L^p(\R^n))$ with $\frac1q+\frac2p=1$, we apply \eqref{eq:holder1} to get
\[
(V_1-V_2)\Psi \in L^{r'}\bigl((0,T);L^{p'}(\R^n)\bigr), \qquad V_2\Phi \in L^{r'}\bigl((0,T);L^{p'}(\R^n)\bigr).
\]

Next, we verify the $L^1(\Sigma)$-bounds. Since $V_j\in L^1(\R^n)\cap L^q(\R^n)$ and $p'\in (1,q)$, we have that $V_1-V_2,\,V_2\in L^{p'}(\R^n)$.  H\"older’s inequality in space yields
\[
\int_{\R^n} |(V_1-V_2)(x)\Psi(t,x)|\,\dd x \le \|V_1-V_2\|_{L^{p'}(\R^n)} \|\Psi(t,\cd)\|_{L^p(\R^n)},
\]
and similarly
\[
\int_{\R^n} |V_2(x)\Phi(t,x)|\,\dd x \le \|V_2\|_{L^{p'}(\R^n)} \|\Phi(t,\cd)\|_{L^p(\R^n)}.
\]
Integrating these inequalities for $t\in(0,T)$ and using that $\Psi,\Phi\in L^{r}\bigl((0,T);L^p(\R^n)\bigr)$, we conclude that
\[
F=(V_1-V_2)\Psi \in L^1(\Sigma), \qquad V_2\Phi \in L^1(\Sigma).
\]
We also record for later use that  $(V_1-V_2)\Psi \ovl{\psi_2}\in L^1(\Sigma)$. This follows from a similar calculation, using $|\psi_2|\le 1+|w_2^{\mathrm{cor}}|$ and $w_2^{\mathrm{cor}}\in L^p(\R^n)$.

Hence, Lemma~\ref{lem:int_by_parts_complex-exp_Lp} applies and yields
\begin{equation}\label{eq:IBP_Phi_exp}
\int_\Sigma e^{i(\lambda^2 t+\lambda\omega_2\cdot x)}(i\partial_t+\Delta)\Phi = 0,
\end{equation}
with $\lambda>0$ and $\omega_2\in\Sph^{n-1}$ as in the statement.

Using that $\Phi$ satisfies \eqref{eq:initial_BVP_F} with $(V_1-V_2)\Psi \ovl{\psi_2}\in L^1(\Sigma)$, we may write
\[
\int_\Sigma (V_1 - V_2)\Psi \ovl{\psi_2}= \int_\Sigma \left(i\d_t + \Delta - V_2\right)\Phi  e^{i\lambda^2 t}\left(e^{i\lambda\omega_2\cdot x}+\ovl{w^{\mathrm{cor}}_2}\right).
\]
Then, expanding the right-hand side and using \eqref{eq:IBP_Phi_exp}, we obtain
\begin{equation}\label{eq:step1}
\int_\Sigma (V_1 - V_2)\Psi \ovl{\psi_2}= - \int_\Sigma \Phi V_2  \ovl{e^{-i(\lambda^2 t+\lambda\omega_2\cdot x)}} + \int_\Sigma  \Phi \ovl{\left(i\d_t+\Delta-\ovl{V_2}\right)\left(e^{-i\lambda^2 t}w^{\mathrm{cor}}_2\right)}.
\end{equation}
However, since $\psi_2=e^{-i\lambda^2 t}(w^{(0)}_2+w^{\mathrm{cor}}_2)$ solves $(i\partial_t+\Delta-\ovl{ V_2})\psi_2=0$, and $e^{-i\lambda^2 t}w^{(0)}_2$ solves the free Schr\"odinger equation, we have
\[
\left(i\d_t+\Delta-\ovl{V_2}\right)\left(e^{-i\lambda^2 t}w^{\mathrm{cor}}_2\right) =\ovl{V_2} e^{-i(\lambda^2 t+\lambda\omega_2\cdot x)}.
\]
The two terms in the right hand side of \eqref{eq:step1} therefore cancel, and we conclude that
\[
\int_\Sigma (V_1 - V_2)\Psi \ovl{\psi_2} = 0.
\]
This completes the proof of the first step.

\vspace{1em}
\noindent\textbf{Step 2:} In this second part of the proof, we use the conclusion established in the first step, applied to solutions $\Psi$ of the Schrödinger equation with potential $V_1$, to deduce that
\[
\int_\Sigma (V_1-V_2)\psi_1\ovl{\psi_2}=0.
\]
To that end, let us consider the final-value problem
\begin{equation}\label{eq:FVP_G}
\begin{cases}
\left(i\d_t + \Delta - \ovl{V_1}\right)\Xi= \left(\ovl{V_1-V_2}\right)\psi_2 \eqqcolon G  \quad &\textrm{in}\quad \Sigma, 
\vspace{.6em}\\
\Xi(T,\cd) = 0   \quad &\textrm{in}\quad \R^n.
\end{cases}
\end{equation}
Using $\eqref{eq:holder1}$ we have $G\in L^{r'}((0,T);L^{p'}(\R^n))$, and hence by the well-posededness discussion in \S\ref{subsec: Stri_WP} there exists a unique solution
\[
\Xi\in C\left([0,T];L^2(\R^n)\right)\cap L^r\left((0,T);L^p(\R^n)\right).
\]
We now verify the assumptions needed to apply Lemma~\ref{lem:CR_lemma4.4_Lp_symmetric} to the solution $\Xi$. As before, we have $ | \psi_2 | \leq 1 + |w^{\mathrm{cor}}_2| $ with $w^{\mathrm{cor}}_2 \in L^p(\R^n)$, and from \eqref{eq:holder1} that $\ovl{G} = (V_1 - V_2)\ovl{\psi_2} \in L^{r'}((0, T); L^{p'}(\R^n))$. Thus $\ovl{G}\in L^{r'}((0,T);L^{p'}(\R^n))$, and since $\Psi\in L^r((0,T);L^p(\R^n))$
the dual pairing $\int_\Sigma \ovl{G} \Psi$ is well defined. Applying the conclusion of the first part of the proof to this $\Psi$, we obtain
\[
\int_\Sigma \ovl{G} \Psi = \int_\Sigma (V_1 - V_2)\Psi\ovl{\psi_2} = 0 .
\]
Using the identity above in Lemma~\ref{lem:CR_lemma4.4_Lp_symmetric} with this choice of $G$ and the solution $\Xi$ of \eqref{eq:FVP_G} yields $\Xi(0,\cd)=0$. Using the assumptions $V_j\in L^1(\R^n)\cap L^q(\R^n)$ and $w^{\mathrm{cor}}_j\in L^p(\R^n)$ for $j\in\{1,2\}$, we have $(V_1-V_2)\psi_1\ovl{\psi_2}\in L^1(\Sigma)$. Moreover, by \eqref{eq:FVP_G} and \eqref{eq:holder}, \eqref{eq:holder1}, 
\[
\left(i\partial_t+\Delta\right)\Xi=\ovl{V_1}\Xi+G \in L^1(\Sigma)\cap L^{r'}\big((0,T);L^{p'}(\R^n)\big).
\]
Since $\Xi$ satisfies~\eqref{eq:FVP_G}, we may write
\[
\int_\Sigma (V_1 - V_2)\psi_1 \ovl{\psi_2} = \int_\Sigma \psi_1 \ovl{\left(i\d_t + \Delta - \ovl{V_1}\right) \Xi}.
\]
Expanding $\psi_1=e^{-i(\lambda^2 t+\lambda\omega_1\cdot x)}+e^{-i\lambda^2 t}w^{\mathrm{cor}}_1$ and replacing in the right hand side of the display above, we obtain
\begin{equation}\label{eq:long}
\begin{split}
\int_\Sigma (V_1 - V_2)\psi_1 \ovl{\psi_2} &= \int_\Sigma e^{-i (\lambda^2 t + \lambda\omega_1\cdot x)} \ovl{\left(i \d_t +\Delta\right ) \Xi} - \int_\Sigma e^{-i (\lambda^2 t + \lambda \omega_1\cdot x)} V_1 \ovl{\Xi} 
\\
&\quad + \int_\Sigma e^{-i\lambda^2 t}w^{\mathrm{cor}}_1 \ovl{\left(i \d_t + \Delta\right) \Xi} - \int_\Sigma e^{-i\lambda^2 t}w^{\mathrm{cor}}_1 V_1 \ovl{\Xi}.
\end{split}
\end{equation}
Now, by Lemma~\ref{lem:int_by_parts_complex-exp_Lp} applied to $\Xi$, we have for the first summand in the right hand side of \eqref{eq:long}
\[
\int_\Sigma e^{-i (\lambda^2 t + \lambda\omega_1\cdot x)} \ovl{\left(i \d_t +\Delta\right) \Xi}=0.
\]
For the third term in the right hand side of \eqref{eq:long}, we apply Proposition~\ref{cor:integration_by_parts_unbounded_no-endopint2}
with $u=\Xi$ and $v=e^{-i\lambda^2 t}w^{\mathrm{cor}}_1$ and obtain
\[
\int_\Sigma e^{-i\lambda^2 t}w^{\mathrm{cor}}_1 \ovl{\left(i \d_t + \Delta\right)\Xi} = \int_\Sigma \ovl{\left(i\d_t+\Delta\right)\left(e^{-i\lambda^2 t}w^{\mathrm{cor}}_1\right)} \Xi.
\]
Since $\psi_1=e^{-i\lambda^2 t}(w^{(0)}_1+w^{\mathrm{cor}}_1)$ solves $(i\partial_t+\Delta-V_1)\psi_1=0$ and $(i\partial_t+\Delta) (e^{-i\lambda^2 t}w^{(0)}_1)=0$,
it follows that
\[
\left(i\d_t+\Delta\right)\left(e^{-i\lambda^2 t}w^{\mathrm{cor}}_1\right) = e^{-i\lambda^2 t}V_1\left(w^{(0)}_1+w^{\mathrm{cor}}_1\right) = V_1 e^{-i (\lambda^2 t + \lambda\omega_1\cdot x)}+ V_1 e^{-i\lambda^2 t}w^{\mathrm{cor}}_1 .
\]
Substituting this identity into the previous display, we get
\[
\int_\Sigma e^{-i\lambda^2 t}w^{\mathrm{cor}}_1 \ovl{\left(i \d_t + \Delta\right)\Xi} = \int_\Sigma V_1 e^{-i \left(\lambda^2 t + \lambda\omega_1\cdot x\right)} \ovl{\Xi} + \int_\Sigma V_1 e^{-i\lambda^2 t}w^{\mathrm{cor}}_1 \ovl{\Xi}.
\]
Therefore the remaining three terms in \eqref{eq:long} cancel, and we conclude that
\[
\int_\Sigma (V_1 - V_2)\psi_1 \ovl{\psi_2}=0,
\]
which proves~\eqref{eq:int_by_parts_time-harmonic}.
\end{proof}
%%%%%%%%%%%%%%%%%%%%%%%%%%%%%% PROOF PROOF PROOF


%%%%%%%%%%%%%%%%%%%%%%%%%%%%%% SECTION SECTION SECTION
\section{The proof of Theorem~\ref{th:Lq}}\label{sec:unbounded_endpoint} We split the proof of Theorem~\ref{th:Lq} into two subsections. The first contains the proof of the non-endpoint case, namely the case that $V\in L^q(\R^n)$ in dimensions $n\geq 2$, for some $q>n/2$. In the subsequent section we will provide the proof of the endpoint case $V\in L^{n/2}(\R^n)$ for dimension $n\geq 3$.

%%%%%%%%%%%%%%%%%%%%%%%%%%%%%% SECTION SECTION SECTION
\subsection{Proof of Theorem~\ref{th:Lq} in the non-endpoint case}\label{subsec:proof_uniqueness_unbounded} Here we consider potentials $V_1,V_2\in L^q(\R^n)$ for some $q>n/2$ and $n\geq 2$. By the discussion in Remark~\ref{rmrk:numer}, we may and do assume that $n/2<q\leq (n+1)/2$, and define $p$ through the relation $\frac{1}{q}=1-\frac{2}{p}$.

 Let $F\coloneqq V_1 - V_2 \in L^1(\R^n)\cap L^q(\R^n)$.  Since $1<p'<q$, we also have that $F\in L^{p'}(\R^n)$. Given $\xi \in \R^n$ consider $\nu \in \Sph^{n-1}$ such that $\xi \cdot \nu = 0$.
For $\lambda \geq |\xi|/2$ we define
\[
\begin{split}
\omega_1 \coloneqq \frac{1}{\lambda} \frac{\xi}{2} + \left( 1 - \frac{|\xi|^2}{4 \lambda^2} \right)^{1/2} \nu,\qquad \omega_2 \coloneqq -\frac{1}{\lambda} \frac{\xi}{2} +\left( 1 - \frac{|\xi|^2}{4 \lambda^2} \right)^{1/2} \nu.
\end{split}
\]
Note that $\omega_1$ and $\omega_2$ belong to $\Sph^{n-1}$ with $\omega_1 - \omega_2 = \xi/\lambda$. We construct the solutions $\psi_1,\psi_2$ solving $(i\partial_t+\Delta-V_1)\psi_1=0$, and $(i\partial_t+\Delta-\ovl{V_2})\psi_2=0$ of the form
\[
\psi_j\coloneqq \psi_{V_j} ^{\lambda,\omega_j}=e^{-i\lambda^2t}w_j(x)=e^{-i\lambda^2t}(w^{(0)}_j(x)+w^{\mathrm{cor}}_j(x)),\qquad j\in\{1,2\},\qquad (t,x)\in\R\times \R^n,
\]
where we remember that $w^{(0)}_j(x)=e^{-i\lambda \omega_j\cdot x}$ and the functions $w^{\mathrm{cor}}_j$ are constructed in Corollary~\ref{cor:solutions} with $\lambda > \max(|\xi|/2, \lambda_{V_j})$. 

In particular, the corrections $w^{\mathrm{cor}}_j$ satisfy the following $L^p$-estimates
\begin{equation}\label{eq:decay_remainder_Lp}
\left\| w^{\mathrm{cor}}_j \right\|_{L^p(\R^n)} \lesssim  \frac{1}{\lambda^{n(\frac{2}{n} - \frac{1}{q})}} \| V_j \|_{L^{p'} (\R^n)},\qquad j\in\{1,2\},
\end{equation}
for all $\lambda \geq \max(|\xi|/2, \lambda_{V_j} + \delta)$. Since $ \mathcal{U}^1_T=\mathcal{U}^2_T$, plugging the solutions $\psi_1,\psi_2$ into the orthogonality relation of Proposition~\ref{pr:int_by_parts_time-harmonic_solutions}, and noting that the time-dependent phases cancel, we infer that
\begin{equation}\label{eq:leading=remainder_unbounded}
\int_{\R^n} F w^{(0)}_1 \ovl{w^{(0)}_2}  = - \int_{\R^n} F \left[w^{(0)}_1 \ovl{w^{\mathrm{cor}}_2} + w^{\mathrm{cor}}_1 \ovl{w^{(0)}_2} + w^{\mathrm{cor}}_1 \ovl{w^{\mathrm{cor}}_2}\right].
\end{equation}
For the left-hand side of \eqref{eq:leading=remainder_unbounded}, note that  
\[
\int_{\R^n} F w^{(0)}_1 \ovl{w^{(0)}_2}  = \int_{\R^n} F(x) e^{-i \lambda( \omega_1-\omega_2)\cdot x}\,\dd x  = \int_{\R^n} F(x) e^{-i \xi \cdot x} \,\dd x = (2 \pi)^{n/2} \wh{F}(\xi).
\]
Using H\"older's inequality along with the estimates \eqref{eq:decay_remainder_Lp}, the right-hand side of \eqref{eq:leading=remainder_unbounded} is readily estimated as follows
\[
\begin{split}
|\wh{F}(\xi)| &\lesssim  \| F \|_{L^{p'}(\R^n)} \left\|w^{\mathrm{cor}}_2\right\|_{L^p(\R^n)} + \left\| F \right\|_{L^{p'}(\R^n)} \left\|w^{\mathrm{cor}}_1\right\|_{L^p(\R^n)} 
\\
&\qquad + \| F \|_{L^q(\R^n)} \left\| w^{\mathrm{cor}}_1 \right\|_{L^p(\R^n)} \left\| w^{\mathrm{cor}}_2 \right\|_{L^p(\R^n)} 
\\
&\lec \frac{1}{\lambda^{n(\frac{2}{n} - \frac{1}{q})}} \| F \|_{L^{p'}(\R^n)} \left(\| V_1 \|_{L^{p'}(\R^n)} +  \| V_2 \|_{L^{p'}(\R^n)} \right) 
\\
&\qquad + \frac{1}{\lambda^{2n(\frac{2}{n} - \frac{1}{q})}} \| F \|_{L^{q}(\R^n)}  \| V_1 \|_{L^{p'}(\R^n)} \| V_2 \|_{L^{p'}(\R^n)} \\
&\lec  \frac{1}{\lambda^{n(\frac{2}{n} - \frac{1}{q})}} + \frac{1}{\lambda^{2n(\frac{2}{n} - \frac{1}{q})}} .
\end{split}
\]
Letting $\lambda\to \infty $ we obtain that $\wh{F}(\xi) = 0$ since $q>n/2$. Since $\xi\in\R^n$ was arbitrary this implies that $V_1=V_2$ almost everywhere in $\R^n$, thus concluding the proof of Theorem~\ref{th:Lq} in the case $q>n/2$ and any dimension $n\geq 2$.


%%%%%%%%%%%%%%%%%%%%%%%%%%%%%% SECTION SECTION SECTION
\subsection{Proof of Theorem~\ref{th:Lq}: the endpoint case}\label{subsec: proof_uniqueness_endpoint} Let $V_1,V_2\in L^1(\R^n)\cap L^{n/2}(\R^n)$ and define $F\coloneqq V_1 - V_2$. Given $\xi \in \R^n$ we consider $\nu \in \Sph^{n-1}$ such that $\xi \cdot \nu = 0$ and $\lambda \geq \max\{\lambda_{V_1},\lambda_{V_2},|\xi|/2\}$, where $\lambda_{V_1},\lambda_{V_2}$ are the constants corresponding to $V_1,V_2$ from Corollary~\ref{cor:solutionsVcrit}. As in \S\ref{subsec:proof_uniqueness_unbounded}, we consider the vectors
\[
\omega_1 \coloneqq \frac{1}{\lambda} \frac{\xi}{2} + \left( 1 - \frac{|\xi|^2}{4 \lambda^2} \right)^{1/2} \nu,\qquad \omega_2 \coloneqq -\frac{1}{\lambda} \frac{\xi}{2} + \left( 1 - \frac{|\xi|^2}{4 \lambda^2} \right)^{1/2} \nu,
\]
and note that $\omega_1,\omega_2\in\mathbb S^{n-1}$ and $\omega_1 - \omega_2 = \xi/\lambda$.  
We now construct the stationary state solutions as in Corollary~\ref{cor:solutionsVcrit}
\[
\psi_j\coloneqq \psi_{V_j} ^{\lambda,\omega_j}=e^{-i\lambda^2t}w_j(x)=e^{-i\lambda^2t}\left(w^{(0)}_j(x)+w^{\mathrm{cor}}_j(x)\right),\qquad j\in\{1,2\},\qquad (t,x)\in\R\times \R^n,
\]
and we recall the estimates
\begin{equation}\label{eq:psiestimates}
\left\|w^{\mathrm{cor}}_j\right\|_{X_\lambda ^*}\lesssim \|V_j\|_{X_\lambda}\leq \lambda^{-\frac{1}{n+1}} \|V_j\|_{L^1(\R^n)\cap L^{n/2}(\R^n)},\qquad j\in\{1,2\},
\end{equation}
also proved in Corollary~\ref{cor:solutionsVcrit}. Since $ \mathcal{U}^1_T=\mathcal{U}^2_T$, plugging the solutions $\psi_1,\psi_2$ into the orthogonality relation of 
Proposition~\ref{pr:int_by_parts_time-harmonic_solutions} and repeating the calculations from \S\ref{subsec:proof_uniqueness_unbounded}, we have 
\[
\begin{split}
|\wh{F}(\xi)| &\leq \frac{1}{(2\pi)^{n/2}} \int_{\R^n} \left|F w^{(0)}_1 \overline{w^{\mathrm{cor}}_2}\right| + \int_{\R^n} \left|F \overline{w^{(0)}_2} w^{\mathrm{cor}}_1\right| + \int_{\R^n} \left|F w^{\mathrm{cor}}_1 \overline{w^{\mathrm{cor}}_2}\right| 
\\
&\lesssim \left\|Fw^{(0)}_1\right\|_{X_\lambda} \left\|w^{\mathrm{cor}}_2\right\|_{\X_\lambda ^*}+\left\|Fw^{(0)}_2\right\|_{X_\lambda} \left\|w^{\mathrm{cor}}_1\right\|_{\X_\lambda ^*}+\|F\|_\lambda\|w^{\mathrm{cor}}_1\|_{\X_\lambda ^*} \left \|w^{\mathrm{cor}}_2\right \|_{\X_\lambda ^*},
\end{split}
\]
where we have used the duality of $X_\lambda$ and $X_\lambda ^*$ and the H\"older-type inequality of Lemma~\ref{lem:multiplication_V_endpoint_X_norm}. Using the estimates \eqref{eq:psiestimates} and the fact that $|w^{(0)}_1|=|w^{(0)}_2|=1$ we get
\begin{align*}
|\wh{F}(\xi)| &\lesssim 
 \|F\|_{\X_\lambda} \lambda^{-\frac{1}{n+1}}(\|V_1\|_{L^1(\R^n)\cap L^{n/2}(\R^n)} + \|V_2\|_{L^1(\R^n)\cap L^{n/2}(\R^n)}) \\
&+ \|F\|_{\lambda} \lambda^{-\frac{2}{n+1}} \|V_1\|_{L^1(\R^n)\cap L^{n/2}(\R^n)} \|V_2\|_{L^1(\R^n)\cap L^{n/2}(\R^n)} \\
%\|F\|_{X_\lambda} ^2\left(1+\|F\|_{\lambda}\right)
&\lesssim\lambda^{-\frac{2}{n+1}} \left(\|V_1\|_{L^1(\R^n)\cap L^{n/2}(\R^n)} +\|V_2\|_{L^1(\R^n)\cap L^{n/2}(\R^n)}\right) ^2 
(1 + \|F\|_{\lambda}) \to 0
\end{align*}
as $\lambda \to \infty$. Hence $\wh{F}(\xi) = 0$, and since $\xi\in\R^n$ was arbitrary we conclude that $F=0$ and hence that $V_1=V_2$ almost everywhere in $\R^n$, thus completing the proof of the endpoint case of Theorem~\ref{th:Lq}.

%%%%%%%%%%%%%%%%%%%%%%%%%%%%%% ACKNOWLEDGMENTS ACKNOWLEDGMENTS ACKNOWLEDGMENTS
\sloppy
\begin{acknowledgements}
M. Ca\~nizares is partially supported by grant PID2024-156267NB-I00 funded by MICIU/AEI/10.13039/501100011033 and cofunded by the European Union. P. Caro is supported by grant PID2024-156267NB-I00 funded by MICIU/AEI/10.13039/501100011033 and cofunded by the European Union, as well as BCAM-BERC 2022-2025 and the BCAM Severo Ochoa CEX2021-001142-S. T. Zacharopoulos is supported by the grant 10.46540/3120-00003B from Independent Research Fund Denmark. I. Parissis is partially supported by grant PID2024-156267NB-I00 funded by MICIU/AEI/10.13039/501100011033 and cofunded by the European Union, grant IT1615-22 of the Basque Government and IKERBASQUE.
\end{acknowledgements}
%%%%%%%%%%%%%%%%%%%%%%%%%%%%%% ACKNOWLEDGMENTS ACKNOWLEDGMENTS ACKNOWLEDGMENTS

%%%%%%%%%%%%%%%%%%%%%%%%%%%%%% BIBLIOGRAPHY BIBLIOGRAPHY BIBLIOGRAPHY
\bibliography{references}{} 
\end{document}

